\documentclass[11pt]{article}
\usepackage[a4paper,margin=23mm,headheight=15pt]{geometry}
\usepackage[T1]{fontenc}
\usepackage{lmodern,microtype}
\usepackage{amsmath,amssymb,amsthm,mathtools}
\usepackage{graphicx,xcolor,booktabs,longtable,array,tabularx}
\usepackage{enumitem,fancyhdr,lastpage,listings,float}
\usepackage{xurl}
\usepackage[hidelinks]{hyperref}
\definecolor{navy}{HTML}{16354A}
\definecolor{teal}{HTML}{007D80}
\definecolor{amber}{HTML}{B57524}
\definecolor{light}{HTML}{F0F5F7}
\definecolor{ink}{HTML}{263846}
\color{ink}
\setlength{\emergencystretch}{3em}
\setlength{\parindent}{0pt}
\setlength{\parskip}{0.5em}
\setlist{itemsep=0.2em,topsep=0.3em}
\newtheorem{theorem}{Theorem}[section]
\newtheorem{proposition}[theorem]{Proposition}
\newtheorem{lemma}[theorem]{Lemma}
\newtheorem{corollary}[theorem]{Corollary}
\theoremstyle{definition}
\newtheorem{definition}[theorem]{Definition}
\newtheorem{example}[theorem]{Example}
\theoremstyle{remark}
\newtheorem{remark}[theorem]{Remark}
\newcommand{\Q}{\mathbb Q}
\newcommand{\R}{\mathbb R}
\newcommand{\N}{\mathbb N}
\newcommand{\E}{\mathbb E}
\newcommand{\Prob}{\mathbb P}
\newcommand{\avocdate}{7 October 2026}
\newcommand{\avocsubtitle}{Certified Artifact Planning, Calibration, and Recoverable Execution}
\providecommand{\paperkind}{Publication edition}
\newcommand{\note}[2]{\begin{quote}\textcolor{teal}{\textbf{#1}}\par #2\end{quote}}
\newcommand{\fig}[3]{\begin{figure}[H]\centering\includegraphics[width=\linewidth]{#1}\caption{#2}\label{#3}\end{figure}}
\graphicspath{{../figures/}{figures/}}
\pagestyle{fancy}\fancyhf{}
\fancyhead[L]{\small\textcolor{navy}{AVOC / Certified artifact planning}}
\fancyhead[R]{\small\paperkind}
\fancyfoot[L]{\small Komorebi P.\quad $\cdot$\quad \avocdate}
\fancyfoot[R]{\small\thepage\,/\,\pageref{LastPage}}
\renewcommand{\headrulewidth}{0.3pt}
\hypersetup{pdfauthor={Komorebi P.},pdftitle={AVOC: Certified Artifact Planning, Calibration, and Recoverable Execution},pdfsubject={Publication edition based on AVOC Formal v0.41.0}}
\lstdefinelanguage{Lean}{keywords={import,namespace,end,open,def,theorem,lemma,structure,inductive,where,by,exact,apply,intro,have,show,match,with,if,then,else,fun,let,in,private,noncomputable,instance,variable,abbrev,example,deriving},sensitive=true,morecomment=[l]{--},morecomment=[s]{/-}{-/},morestring=[b]"}
\lstdefinestyle{LeanCode}{language=Lean,basicstyle=\ttfamily\footnotesize,breaklines=true,columns=fullflexible,keepspaces=true,showstringspaces=false,backgroundcolor=\color{light},frame=single,rulecolor=\color{light},keywordstyle=\color{navy},commentstyle=\color{teal},aboveskip=0.7em,belowskip=0.7em,
 literate={∀}{{$\forall$}}1 {∃}{{$\exists$}}1 {→}{{$\to$}}1 {↔}{{$\leftrightarrow$}}1 {⇒}{{$\Rightarrow$}}1 {←}{{$\leftarrow$}}1 {∈}{{$\in$}}1 {∉}{{$\notin$}}1 {≤}{{$\le$}}1 {≥}{{$\ge$}}1 {≠}{{$\ne$}}1 {ℕ}{{$\mathbb N$}}1 {ℚ}{{$\mathbb Q$}}1 {ℝ}{{$\mathbb R$}}1 {ℤ}{{$\mathbb Z$}}1 {Ω}{{$\Omega$}}1 {α}{{$\alpha$}}1 {β}{{$\beta$}}1 {κ}{{$\kappa$}}1 {δ}{{$\delta$}}1 {ε}{{$\varepsilon$}}1 {π}{{$\pi$}}1 {ω}{{$\omega$}}1 {σ}{{$\sigma$}}1 {τ}{{$\tau$}}1 {μ}{{$\mu$}}1 {λ}{{$\lambda$}}1 {Θ}{{$\Theta$}}1 {ξ}{{$\xi$}}1 {η}{{$\eta$}}1 {ν}{{$\nu$}}1 {χ}{{$\chi$}}1 {θ}{{$\theta$}}1 {ρ}{{$\rho$}}1 {Γ}{{$\Gamma$}}1 {Σ}{{$\Sigma$}}1 {∧}{{$\land$}}1 {∨}{{$\lor$}}1 {¬}{{$\neg$}}1 {×}{{$\times$}}1 {⊆}{{$\subseteq$}}1 {⊂}{{$\subset$}}1 {⟨}{{$\langle$}}1 {⟩}{{$\rangle$}}1 {∑}{{$\sum$}}1 {∏}{{$\prod$}}1 {∣}{{$\mid$}}1 {∘}{{$\circ$}}1 {↑}{{$\uparrow$}}1 {↓}{{$\downarrow$}}1 {▸}{{$\triangleright$}}1 {₀}{{$_0$}}1 {₁}{{$_1$}}1 {₂}{{$_2$}}1 {₃}{{$_3$}}1 {₄}{{$_4$}}1 {₅}{{$_5$}}1 {₆}{{$_6$}}1 {₇}{{$_7$}}1 {₈}{{$_8$}}1 {₉}{{$_9$}}1}
\lstset{style=LeanCode}
\newcommand{\avoccover}[2]{\hypersetup{pageanchor=false}\begin{titlepage}\color{navy}
 \vspace*{16mm}{\large KOMOREBI\quad /\quad REALIZATION CALCULUS\par}
 \vspace{17mm}{\Huge\bfseries AVOC\par}
 \vspace{9mm}{\LARGE\bfseries Certified Artifact Planning,\\[2mm]Calibration, and\\[2mm]Recoverable Execution\par}
 \vspace{14mm}{\Large #1\par}\vspace{10mm}\rule{\linewidth}{0.6pt}
 \vspace{8mm}{\normalsize #2\par}
 \vfill{\large Komorebi P.\par}\vspace{3mm}{\avocdate\quad /\quad Publication edition 1.0\par}
 \end{titlepage}\hypersetup{pageanchor=true}}

\usepackage{xurl}
\lstdefinestyle{AVOCSource}{style=LeanCode,
  literate={ℚ}{{$\mathbb Q$}}1 {ℕ}{{$\mathbb N$}}1 {ℝ}{{$\mathbb R$}}1
    {ℤ}{{$\mathbb Z$}}1 {Ω}{{$\Omega$}}1 {ω}{{$\omega$}}1
    {α}{{$\alpha$}}1 {β}{{$\beta$}}1 {δ}{{$\delta$}}1
    {ε}{{$\varepsilon$}}1 {η}{{$\eta$}}1 {ζ}{{$\zeta$}}1
    {φ}{{$\varphi$}}1 {ψ}{{$\psi$}}1 {χ}{{$\chi$}}1
    {∀}{{$\forall$}}1 {∃}{{$\exists$}}1 {→}{{$\to$}}1
    {←}{{$\leftarrow$}}1 {↔}{{$\leftrightarrow$}}1
    {∧}{{$\land$}}1 {∨}{{$\lor$}}1 {¬}{{$\neg$}}1
    {∈}{{$\in$}}1 {∉}{{$\notin$}}1 {∅}{{$\varnothing$}}1
    {⊆}{{$\subseteq$}}1 {≤}{{$\le$}}1 {≠}{{$\ne$}}1
    {∑}{{$\sum$}}1 {∏}{{$\prod$}}1 {×}{{$\times$}}1
    {∘}{{$\circ$}}1 {⟨}{{$\langle$}}1 {⟩}{{$\rangle$}}1
    {↑}{{$\uparrow$}}1 {↓}{{$\downarrow$}}1 {▸}{{$\triangleright$}}1
    {₀}{{$_0$}}1 {₁}{{$_1$}}1 {₂}{{$_2$}}1 {₃}{{$_3$}}1 {₄}{{$_4$}}1
    {¹}{{$^1$}}1 {⁻}{{$^-$}}1 {‹}{{$\langle$}}1 {›}{{$\rangle$}}1
    {«}{{\guillemotleft}}1 {»}{{\guillemotright}}1 {·}{{$\cdot$}}1 {⊢}{{$\vdash$}}1}
\title{AVOC - Lean Companion\\[0.5em]\large Certified Artifact Planning, Calibration, and Recoverable Execution}
\author{Komorebi P.}
\date{\avocdate}

% Operation Komorebi: shared publication profile.
\usepackage{komorebi-publication}
\renewcommand{\KomorebiTitle}{AVOC}
\renewcommand{\KomorebiSubtitle}{Artifact value, guarded operations, and calibrated planning}
\renewcommand{\KomorebiSeries}{Research library}
\renewcommand{\KomorebiKind}{Lean and evidence guide}
\renewcommand{\KomorebiEdition}{Standardized publication edition / 8 October 2026}
\renewcommand{\KomorebiShortTitle}{AVOC}
\renewcommand{\KomorebiTitleSize}{38}
\renewcommand{\KomorebiTitleLeading}{43}
\renewcommand{\KomorebiProfile}{classical}
\renewcommand{\KomorebiOrientation}{portrait}
\KomorebiSetup
\begin{document}
\KomorebiCover
\KomorebiOpeningNote{Actual source declarations, proof structure, module maps, negative controls, and pinned reproduction instructions. Archived kernel evidence is distinguished from fresh source and receipt checks.}

\renewcommand{\paperkind}{Lean companion}

\begin{abstract}
This companion explains the cumulative AVOC v0.41.0 formalization through its
actual source declarations, hypotheses, proof structure, and executable
contracts. It connects information refinement and finite decision theory to
guarded planning, semantic artifact verification, cost accounting, recovery,
finite calibration, and operation lifecycle recording. The final statistical
composition bounds the unconditional probability of issuing an uncertified
policy episode under a finite rational source model. Source-faithful excerpts,
module maps, negative controls, and reproduction commands make each claim
traceable. The retained release records report 222 clean local target builds,
2,268 named theorem/lemma axiom reports, and 135 substantive negative-test
rejections. This publication performs fresh source and receipt checks but does
not claim a fresh Lean compilation.
\end{abstract}
\tableofcontents

\section{How to read the formal evidence}
\label{sec:evidence}
The main paper states the mathematical and operational results. The visual
guide follows concrete examples. This companion answers a different question:
which precise Lean objects and declarations support each result, and what must
a user supply before the theorem applies? It is written for a mathematician or
developer who wants to inspect that boundary without reading every module in
sequence.

The formal checkpoint is \path{AVOC_Formal_v0.41.0}. Its umbrella
\path{AVOC.lean} imports the cumulative theory. Its source archive retains
earlier development failures, negative controls, independent reviews, runtime
fixtures, and historical release records. The publication does not edit those
sources. An assertion about a theorem below refers to that checkpoint, rather
than to a newly created proof file.

\subsection{Three kinds of evidence}
There are three separate objects to inspect. A Lean theorem is a universally
quantified statement about its defined objects, with explicit hypotheses. A
kernel audit establishes that the elaborated proof term depends only on the
reported axioms. A finite executable comparison establishes agreement on the
declared tested family. None of these should be silently substituted for
another.

For example, the theorem that recorder samples form a source prefix quantifies
over every successful finite run of the defined generative machine. The
69,905-word Python comparison exercises a bounded family of raw recorder
commands. The latter provides adapter evidence. It does not prove the
generative theorem or identify a physical provider with its source model.

\begin{center}
\begin{tabular}{p{0.28\linewidth}p{0.23\linewidth}p{0.38\linewidth}}
\toprule
Object & Reported scope & Interpretation\\
\midrule
Mathematical modules & 195 non-audit modules & Source organization, including modules with no named theorem\\
Named theorem/lemma inventory & 2,268 declarations & Includes two private helpers, and many elementary supporting results\\
Clean local builds & 222 targets & Mathematical modules, umbrella and executable/reporting targets\\
Kernel axiom reports & 2,268 exact names & Archived matching membership, with no missing or duplicate report\\
Negative controls & 135 substantive rejections & Archived false proof attempts rejected for mathematical reasons\\
Fresh publication check & Source and receipt consistency & No kernel or Lean endpoint execution here\\
\bottomrule
\end{tabular}
\end{center}

These counts describe verification coverage. They do not count independent
discoveries. The inventory is more useful as a navigable list of exact
declarations than as a single headline number.

\subsection{Compiler and axiom provenance}
The pinned toolchain is Lean 4.34.0-rc2, commit
\path{6a10ac8c22beadecabdbb0919c2b50214762f91d}. The pinned mathlib commit is
\path{aeb72f1e9981b6d35b426b16dca7759be3e7ba06}. The release's union of reported
axioms is \path{propext}, \path{Classical.choice}, and \path{Quot.sound}.
This is the standard classical Lean foundation used by the checkpoint.
The recorded clean builds used warnings as errors and rebuilt local objects
from empty. External dependency caches were checked against pins. Their entire
source trees were not rebuilt as part of that local build claim.

The present environment has no available Lean or Lake compiler. Therefore
\textbf{fresh kernel compilation is not run here}. The companion preserves
the distinction in its inventory JSON, captions, and reproduction instructions.
It supplies no newly invented axiom report and no new proof file presented as
compiled.

\section{A map of the cumulative theory}
AVOC starts with representations and targets, then adds costs, evidence,
actions, probability, and persistent protocol state. The later recorder block
does not replace the earlier theory. It closes a particular lifecycle
correspondence and composes the actual resulting samples with the existing
calibration and policy statements.

\begin{center}
\begin{tabular}{p{0.24\linewidth}p{0.34\linewidth}p{0.33\linewidth}}
\toprule
Layer & Representative modules & Main obligation\\
\midrule
Information & Core, Refinement, RepresentationLattice & A representation preserves the declared target\\
Decision & Control, Dynamic, Stochastic, Bayesian & An action or plan has its stated value\\
Approximation & ApproximateStochastic, QRateDistortion, GuardedModelError & Quantified error propagates to policy regret\\
Artifacts & CertificateLogic, ArtifactBinding, BooleanCostClosure, TypedAffineProgram & Submitted bytes and semantics satisfy a specific checker contract\\
Operational control & AttemptAccounting, VerifiedUpgrade, LazyStagedControl & Paid actions preserve credit and legal state\\
Recovery & CommitRecovery, DurableUpgrade, RobustQuoteRecovery & Recorded committed transitions recover safely\\
Uncertain quotes & QuoteAuthority, RobustQuotePlanning & One compatible realization supports intervals and regret\\
Calibration & FiniteCalibration, CalibrationSchedule, ObservationStream & A fixed scheduled interval construction has its coverage bound\\
Observation lifecycle & ObservationRecorder, RecorderReveal, RecorderPolicy & Actual model-generated terminal samples feed the issued policy\\
\bottomrule
\end{tabular}
\end{center}

This table is a reading map, not a complete import graph. The source inventory
records every immediate import. A useful distinction is between an abstract
existence theorem, an executable implementation theorem, and a proof that a
particular operational model realizes the abstract object. The program often
contains all three, but in different files and at different levels of scope.

\subsection{Lean conventions used below}
\path{Prop} contains propositions, while \path{Type} contains data. A
\path{structure} can hold both data fields and proof fields. For example, a
rational probability row carries nonnegativity and a sum-to-one proof.
\path{Finset} is a duplicate-free finite set. \path{List} is an ordered sequence,
where duplicates must be controlled separately. \path{Fin (n+1)} represents an
integer count between zero and $n$, with that bound carried in its type.

\path{Option} distinguishes an accepted result from no result. \path{Except}
additionally retains an error value. A \path{noncomputable def} may use a
classical existence selector. It establishes a mathematical value or policy,
not an executable runtime bound. The executable planner and codec modules are
the place to inspect algorithms and decoding.

Source excerpts below preserve the actual Lean spelling. Mathematical symbols
are rendered through the listings style. A statement excerpt stops before its
proof body and is not presented as a standalone Lean file. Full proof excerpts
are explicitly identified. The complete source remains the authority for
surrounding variables, instances, and imports.

\section{Representations, fibers, and exact sufficiency}
\label{sec:fibers}
The foundational preorder is deterministic factorization. If
$\phi:X\to A$ refines $\psi:X\to B$, there is a decoder $d:A\to B$ such that
$d\circ\phi=\psi$. No particular byte format, compression algorithm, or
probability law is assumed here. The declared target $T:X\to Y$ is sufficient
from $\phi$ exactly when it factors through $\phi$.

\paragraph{Full proof excerpt.}\path{Komorebi.AVOC.refines_trans}. Source: \path{AVOC/Core.lean}, line 32.
\begin{lstlisting}[style=AVOCSource]
@[trans] theorem refines_trans {X : Type u} {A : Type v} {B : Type w} {C : Type z}
    {φ : X → A} {ψ : X → B} {χ : X → C}
    (hφψ : Refines φ ψ) (hψχ : Refines ψ χ) : Refines φ χ := by
  rcases hφψ with ⟨d₁, hd₁⟩
  rcases hψχ with ⟨d₂, hd₂⟩
  refine ⟨d₂ ∘ d₁, ?_⟩
  intro x
  simp only [Function.comp_apply, hd₁ x, hd₂ x]
\end{lstlisting}


The proof builds the composite decoder. Transitivity is not an informal
analogy between encodings. It is function composition with the two
commuting-square equations inserted into the proof term.

\subsection{The collision obstruction}
The simplest failure certificate is a pair of inputs that the representation
merges while the target distinguishes them.

\paragraph{Full proof excerpt.}\path{Komorebi.AVOC.collision_obstructs_sufficiency}. Source: \path{AVOC/Core.lean}, line 87.
\begin{lstlisting}[style=AVOCSource]
theorem collision_obstructs_sufficiency
    {X : Type u} {Z : Type v} {Y : Type w}
    {φ : X → Z} {T : X → Y} {x y : X}
    (hφ : φ x = φ y) (hT : T x ≠ T y) : ¬ TaskSufficient φ T := by
  intro hs
  exact hT (sufficient_preserves_target_distinctions hs hφ)
\end{lstlisting}


The proof uses the supplied decoder to derive equality of the two target
values, then contradicts the target distinction. A reviewer can inspect this
before considering any cardinality or optimization statement. One genuine
task-relevant collision already rules out exact sufficiency.

\paragraph{Statement excerpt.}\path{Komorebi.AVOC.sufficient_iff_fiber_refinement}. Source: \path{AVOC/Core.lean}, line 124.
\begin{lstlisting}[style=AVOCSource]
theorem sufficient_iff_fiber_refinement
    {X : Type u} {Z : Type v} {Y : Type w} [Nonempty Y]
    {φ : X → Z} {T : X → Y} :
    TaskSufficient φ T ↔ ∀ x y, SameFiber φ x y → SameFiber T x y
\end{lstlisting}


The reverse implication needs \path{Nonempty Y}, because a decoder must be
defined even on representation states never reached by an input. Lean uses
classical choice for a representative of a reached fiber and a default value
elsewhere. The condition is mathematical, not a proof that an efficient
decoder can be recovered from arbitrary input data.

For a family of queries, \path{JointTarget} stores the complete vector of
answers. \path{sufficient_joint_iff_each} shows that exact recovery of every
query is equivalent to exact recovery of that vector. The target is a
canonical coarsest sufficient representation up to mutual factorization.
This uniqueness is relative to the declared query family. It does not make a
representation adequate for future, undeclared questions.

\section{Costs turn information into a decision problem}
Information preservation alone cannot choose whether to compute. The
\path{Control} module supplies action scores, feasible metrics, a stop action,
confidence-sensitive answering, and a weak-link generation/verification model.
Its finite maximizer results use a nonempty feasible action set. They do not
claim that finding the maximizer is free.

\paragraph{Full proof excerpt.}\path{Komorebi.AVOC.weakLinkUtility_positive_iff}. Source: \path{AVOC/Control.lean}, line 174.
\begin{lstlisting}[style=AVOCSource]
theorem weakLinkUtility_positive_iff
    (p value generationCost verificationCost : ℝ) :
    0 < weakLinkUtility p value generationCost verificationCost ↔
      generationCost + verificationCost < p * value := by
  unfold weakLinkUtility
  constructor <;> intro h <;> nlinarith
\end{lstlisting}


This arithmetic threshold states exactly when expected admitted value exceeds
generation plus verification cost. It does not prove the input success
probability, monetary value, or cost measurement. Those quantities are model
inputs. Later quote authority and calibration modules explain how some inputs
can be bound to retained evidence.

\path{stop_optimal_of_nonpositive} formalizes a genuine option to stop. It
prevents a model from being forced into a negative-value action merely because
all available computations are unattractive. The queue lemmas separately
track the imbalance between generation and absorption. Increasing generation
capacity can increase an unverified backlog without improving admitted yield.

\section{Finite stochastic control and exact quotients}
The finite stochastic layer uses rational probability tables. Its
\path{FiniteStochasticControlSystem} carries an admissible action set, a
transition row for every state/action pair, row nonnegativity, row normalization,
and rational rewards. Expected continuation is a finite weighted sum.

\paragraph{Statement excerpt.}\path{Komorebi.AVOC.stochastic_finite_horizon_optimal_action_exists}. Source: \path{AVOC/Stochastic.lean}, line 81.
\begin{lstlisting}[style=AVOCSource]
theorem stochastic_finite_horizon_optimal_action_exists
    {S : Type u} {A : Type v} [Fintype S] [DecidableEq A]
    (sys : FiniteStochasticControlSystem S A) (n : ℕ) (s : S) :
    ∃ a ∈ sys.actions,
      stochasticHorizonValue sys (n + 1) s = stochasticQValue sys n s a
\end{lstlisting}


The Bellman recurrence has an optimal action at every finite horizon. This
result is an existence and equality statement. The compiler's treatment of the
definition as \path{noncomputable} matters if an application needs executable
search rather than a proof that a maximum exists.

\subsection{What an exact stochastic quotient must preserve}
A deterministic encoding can preserve reward yet fail to preserve future
control. Exact stochastic abstraction requires the relevant aggregated
transition probabilities to be constant on fibers as well. The quotient proof
groups the raw expectation by coarse state, then transports the Bellman
recurrence through that equality.

\paragraph{Statement excerpt.}\path{Komorebi.AVOC.quotientStochasticControlSystem_horizonValue_eq}. Source: \path{AVOC/StochasticQuotient.lean}, line 151.
\begin{lstlisting}[style=AVOCSource]
theorem quotientStochasticControlSystem_horizonValue_eq
    {S : Type u} {A : Type v} {Z : Type w}
    [Fintype S] [Fintype Z] [DecidableEq Z] [DecidableEq A]
    (sys : FiniteStochasticControlSystem S A)
    (φ : S → Z) (hsurj : Function.Surjective φ)
    (h : ProjectedStochasticControlSufficient sys φ) :
    ∀ n s,
      stochasticHorizonValue (quotientStochasticControlSystem sys φ hsurj) n (φ s) =
        stochasticHorizonValue sys n s
\end{lstlisting}

\paragraph{Statement excerpt.}\path{Komorebi.AVOC.quotient_stochasticBellmanOptimalAt_iff}. Source: \path{AVOC/StochasticQuotient.lean}, line 219.
\begin{lstlisting}[style=AVOCSource]
theorem quotient_stochasticBellmanOptimalAt_iff
    {S : Type u} {A : Type v} {Z : Type w}
    [Fintype S] [Fintype Z] [DecidableEq Z] [DecidableEq A]
    (sys : FiniteStochasticControlSystem S A)
    (φ : S → Z) (hsurj : Function.Surjective φ)
    (h : ProjectedStochasticControlSufficient sys φ)
    (n : ℕ) (s : S) (a : A) :
    StochasticBellmanOptimalAt (quotientStochasticControlSystem sys φ hsurj) n (φ s) a ↔
      StochasticBellmanOptimalAt sys n s a
\end{lstlisting}


These declarations identify value and optimality in the source and quotient
systems under the stated sufficient-abstraction hypotheses. They are not
claims about unrestricted stochastic differential equations. Their scope is
the finite rational control objects named in the signatures.

\section{Approximate abstraction has an explicit horizon cost}
An approximate stochastic abstraction supplies reward error $\varepsilon$,
transition-row $\ell^1$ error $\delta$, and a reward envelope $R$. The proof
separates pointwise value error from transport of the continuation value.
Nonnegative normalized weights keep pointwise error unchanged under
expectation. Row error multiplies the remaining value envelope.

\paragraph{Statement excerpt.}\path{Komorebi.AVOC.approxBudget_closed_form}. Source: \path{AVOC/ApproximateStochastic.lean}, line 264.
\begin{lstlisting}[style=AVOCSource]
theorem approxBudget_closed_form (ε δ R : ℚ) (n : ℕ) :
    approxBudget ε δ R n =
      (n : ℚ) * ε + δ * R * ((n : ℚ) * ((n : ℚ) - 1) / 2)
\end{lstlisting}


The closed form is
\[
 B_n=n\varepsilon+\delta R\frac{n(n-1)}2.
\]
The linear term comes from one reward discrepancy per stage. The triangular
term comes from the remaining-horizon magnitude accumulated through transition
transport. This formula is a finite-horizon bound, not a contraction estimate
or a claim that long-horizon error remains bounded.

\paragraph{Statement excerpt.}\path{Komorebi.AVOC.approximate_stochastic_horizon_value_bound_closed}. Source: \path{AVOC/ApproximateStochastic.lean}, line 386.
\begin{lstlisting}[style=AVOCSource]
theorem approximate_stochastic_horizon_value_bound_closed
    {S : Type u} {A : Type v} {Z : Type w}
    [Fintype S] [Fintype Z] [DecidableEq Z] [DecidableEq A]
    (raw : FiniteStochasticControlSystem S A)
    (absSys : FiniteStochasticControlSystem Z A)
    (φ : S → Z) (ε δ R : ℚ)
    (h : ApproximateStochasticAbstraction raw absSys φ ε δ R)
    (n : ℕ) (s : S) :
    |stochasticHorizonValue raw n s - stochasticHorizonValue absSys n (φ s)| ≤
      (n : ℚ) * ε + δ * R * ((n : ℚ) * ((n : ℚ) - 1) / 2)
\end{lstlisting}


\path{ApproximateComposition} proves transport through successive abstraction
maps, including $\ell^1$ contraction under pushforward. A composed model must
still satisfy the complete combined certificate. Multiplying or adding
informal error summaries is not a replacement for its hypotheses.

\section{Bayesian normalization and its stability boundary}
AVOC's \path{RationalBelief} is a finite normalized nonnegative rational mass
function. Observation channels and evidence likelihoods supply the inputs to
predictive probability and posterior updates. The positive predictive-mass
hypothesis is explicit in the genuine normalized update. A separate totalized
definition handles observations of zero mass, with their contribution to
expected policy value proved to be zero.

\paragraph{Statement excerpt.}\path{Komorebi.AVOC.bayesUpdate_mul_predictive}. Source: \path{AVOC/Bayesian.lean}, line 92.
\begin{lstlisting}[style=AVOCSource]
theorem bayesUpdate_mul_predictive
    {Ω : Type u} [Fintype Ω]
    {O : Type v} [Fintype O]
    (b : RationalBelief Ω) (E : RationalChannel Ω O) (o : O)
    (hpos : 0 < predictiveProb b E o) (ω : Ω) :
    (bayesUpdate b E o hpos).prob ω * predictiveProb b E o =
      b.prob ω * E.prob ω o
\end{lstlisting}

\paragraph{Statement excerpt.}\path{Komorebi.AVOC.bayesianPolicyValue_eq_stochasticDecisionValue}. Source: \path{AVOC/Bayesian.lean}, line 215.
\begin{lstlisting}[style=AVOCSource]
theorem bayesianPolicyValue_eq_stochasticDecisionValue
    {Ω : Type u} [Fintype Ω]
    {O : Type v} [Fintype O]
    {A : Type w}
    (b : RationalBelief Ω) (utility : Ω → A → ℚ)
    (E : RationalChannel Ω O) (policy : O → A) :
    bayesianPolicyValue b utility E policy =
      stochasticDecisionValue b.prob utility E policy
\end{lstlisting}


The first equation lets a proof cancel the normalizer safely. The second
connects signal-conditioned Bayesian decisions with the finite stochastic
decision-value formulation. The equivalence is about those defined values,
not about a black-box posterior engine.

\subsection{Small evidence mass amplifies perturbations}
Approximate conditioning requires a positive lower bound $\zeta$ on the
relevant evidence masses. With prior $\ell^1$ error at most $\alpha$, likelihood
error at most $\beta$, and an upper likelihood bound $K$, the combined theorem
gives a bound of the form
\[
 \|b^L-c^M\|_1\le \frac{2\beta}{\zeta}
                      +\frac{2K\alpha}{\zeta}.
\]

\paragraph{Statement excerpt.}\path{Komorebi.AVOC.conditionByWeight_l1_stability_combined}. Source: \path{AVOC/ApproximateBayesian.lean}, line 227.
\begin{lstlisting}[style=AVOCSource]
theorem conditionByWeight_l1_stability_combined
    {Ω : Type u} [Fintype Ω]
    (b c : RationalBelief Ω)
    (L M : Ω → ℚ)
    (hL : ∀ ω, 0 ≤ L ω) (hM : ∀ ω, 0 ≤ M ω)
    (K α β ζ : ℚ)
    (hK : 0 ≤ K) (hMK : ∀ ω, M ω ≤ K)
    (hα : 0 ≤ α) (hprior : beliefL1Distance b c ≤ α)
    (hβ : 0 ≤ β) (hlik : ∀ ω, |L ω - M ω| ≤ β)
    (hζ : 0 < ζ)
    (hbL : ζ ≤ evidenceMass b L)
    (hbM : ζ ≤ evidenceMass b M)
    (hcM : ζ ≤ evidenceMass c M) :
    beliefL1Distance
        (conditionByWeight b L hL (lt_of_lt_of_le hζ hbL))
        (conditionByWeight c M hM (lt_of_lt_of_le hζ hcM))
      ≤ 2 * β / ζ + 2 * (K * α) / ζ
\end{lstlisting}


The source includes a separate \path{BayesianInstability} module. A rare
observation can sharply amplify differences. The denominator assumptions above
are a substantive boundary, not an inconvenient proof artifact to omit from
the statement. They also explain why compression guarantees must be tailored
to evidence and decision semantics.

\section{Decision compression uses Q signatures}
A representation need not retain every feature of a state to preserve a
particular finite-horizon decision. A Q signature records action values. The
\path{QRateDistortion} module characterizes approximate decoding through fiber
diameters and gives separation lower bounds on the number of codes.

\paragraph{Statement excerpt.}\path{Komorebi.AVOC.exists_qSignatureApproximation_iff_fiberDiameter}. Source: \path{AVOC/QRateDistortion.lean}, line 125.
\begin{lstlisting}[style=AVOCSource]
theorem exists_qSignatureApproximation_iff_fiberDiameter
    {S : Type u} [Fintype S]
    {A : Type v} {Z : Type w} [DecidableEq Z]
    (φ : S → Z) (Q : S → A → ℚ) (η : ℚ) :
    (∃ Qbar : Z → A → ℚ, QSignatureApproximation φ Q Qbar η) ↔
      QFiberDiameterBound φ Q η
\end{lstlisting}

\paragraph{Statement excerpt.}\path{Komorebi.AVOC.q_margin_preserves_unique_optimum}. Source: \path{AVOC/QRateDistortion.lean}, line 256.
\begin{lstlisting}[style=AVOCSource]
theorem q_margin_preserves_unique_optimum
    {S : Type u} {A : Type v} {Z : Type w}
    (φ : S → Z) (Q : S → A → ℚ) (Qbar : Z → A → ℚ) (η : ℚ)
    (happrox : QSignatureApproximation φ Q Qbar η)
    {s : S} {aStar : A}
    (hmargin : ∀ a, a ≠ aStar → 2 * η < Q s aStar - Q s a) :
    ∀ a, a ≠ aStar → Qbar (φ s) a < Qbar (φ s) aStar
\end{lstlisting}


The first result formalizes the radius-versus-diameter relationship through a
midpoint decoder on each finite fiber. The second makes the action margin
explicit: if the original unique maximizer is separated by more than
$2\eta$, approximate decoding with error $\eta$ preserves it. A representation can
therefore be adequate at one horizon or margin and inadequate at another.

The finite code-count bounds are relative to a separated family and the
specified Q coordinates. They do not establish a general theory of arbitrary
neural latent representations. They provide a precise certificate format an
application can try to satisfy.

\section{Guarded plans preserve feasibility on every branch}
The guarded system changes the action set with the state. Its state space need
not be finite, although the observation type and each feasible action set are
finite. A belief state can therefore occur without forcing a finite enumeration
of all possible beliefs.

\texttt{Plan A O n} is an adaptive observation tree of depth $n$. The
\path{Legal} predicate checks the selected action and every continuation,
including branches of zero probability. A comparator may use a different
future action on every observation history. The Bellman optimum still bounds
every such legal finite tree.

\paragraph{Statement excerpt.}\path{Komorebi.AVOC.Guarded.safe_optimal_plan_exists}. Source: \path{AVOC/GuardedControl.lean}, line 177.
\begin{lstlisting}[style=AVOCSource]
theorem safe_optimal_plan_exists (sys : System S A O) (safe : S → Prop)
    (hpres : Preserves sys safe) (n : ℕ) (s : S) (hs : safe s) :
    ∃ p : Plan A O n,
      Legal sys s p ∧ TreeSafe sys safe s p ∧
      planValue sys s p = value sys n s ∧
      ∀ q : Plan A O n, Legal sys s q → planValue sys s q ≤ planValue sys s p
\end{lstlisting}


Safety follows from a proved one-step preservation property for the concrete
system. The declaration does not assume that any external action preserves an
arbitrary invariant. Applications must discharge that obligation with their
actual action semantics.

\section{From local action gaps to whole-policy regret}
\label{sec:regret}
This proof is one of the main reusable connections in the archive. If a policy
chooses an admissible action whose Bellman gap at stage $k$ is at most $e_k$,
the whole expected plan loss is at most $\sum_{k<n}e_k$. The proof is induction
on the horizon and retains every branch's bound.

\paragraph{Full proof excerpt.}\path{Komorebi.AVOC.Guarded.plan_regret_le}. Source: \path{AVOC/GuardedPolicyRegret.lean}, line 25.
\begin{lstlisting}[style=AVOCSource]
theorem plan_regret_le (sys : System S A O) (err : ℕ → ℚ)
    {n : ℕ} (s : S) (p : Plan A O n) (hp : NearGreedy sys err s p) :
    value sys n s - planValue sys s p ≤ regretBudget err n := by
  induction n generalizing s with
  | zero => cases p; simp [value, planValue]
  | succ n ih =>
      cases p with
      | node a next =>
          have hsum : expectation sys (value sys n) s a -
              (∑ o, sys.prob s a o * planValue sys (sys.step s a o) (next o)) ≤
              regretBudget err n := by
            calc
              expectation sys (value sys n) s a -
                  (∑ o, sys.prob s a o * planValue sys (sys.step s a o) (next o))
                = ∑ o, sys.prob s a o *
                    (value sys n (sys.step s a o) -
                     planValue sys (sys.step s a o) (next o)) := by
                       simp [expectation, mul_sub, Finset.sum_sub_distrib]
              _ ≤ ∑ o, sys.prob s a o * regretBudget err n := by
                apply Finset.sum_le_sum
                intro o _
                exact mul_le_mul_of_nonneg_left (ih _ _ (hp.2 o)) (sys.prob_nonneg s a o)
              _ = regretBudget err n := by
                rw [← Finset.sum_mul, sys.prob_sum_one]; simp
          have hlocal := hp.1
          unfold qValue at hlocal
          rw [regretBudget_succ]
          change value sys (n + 1) s -
            (sys.reward s a + ∑ o, sys.prob s a o *
              planValue sys (sys.step s a o) (next o)) ≤ _
          linarith
\end{lstlisting}


At the successor step, the proof writes the difference of expectations as the
expectation of differences. The induction hypothesis bounds every continuation
by the previous regret budget. Nonnegative probabilities permit multiplication
of each inequality. Their sum is one, so that common bound survives averaging.
Adding the local action gap gives the next budget.

\subsection{Approximate Q selection}
Suppose $\widehat Q_k$ approximates each feasible action value within
$\eta_k$, while the selected action is within $\tau_k$ of the estimated
maximum. Comparing the true optimum to the estimate, then to the selected
estimate, then back to the true selected value gives
$2\eta_k+\tau_k$.

\paragraph{Full proof excerpt.}\path{Komorebi.AVOC.Guarded.decoded_action_gap}. Source: \path{AVOC/GuardedPolicyRegret.lean}, line 74.
\begin{lstlisting}[style=AVOCSource]
theorem decoded_action_gap (sys : System S A O) (n : ℕ) (s : S)
    (estimate : A → ℚ) (eta tau : ℚ) (a : A) (ha : a ∈ sys.actions s)
    (he : ∀ b ∈ sys.actions s, |qValue sys n s b - estimate b| ≤ eta)
    (hm : ∀ b ∈ sys.actions s, estimate b ≤ estimate a + tau) :
    value sys (n + 1) s - qValue sys n s a ≤ 2 * eta + tau := by
  have hx := optimalAction_mem sys n s
  have hl := (abs_le.mp (he a ha)).1
  have hu := (abs_le.mp (he (optimalAction sys n s) hx)).2
  have hmax := hm (optimalAction sys n s) hx
  rw [optimalAction_realizes]
  linarith
\end{lstlisting}

\paragraph{Statement excerpt.}\path{Komorebi.AVOC.Guarded.decoded_policy_regret}. Source: \path{AVOC/GuardedPolicyRegret.lean}, line 102.
\begin{lstlisting}[style=AVOCSource]
theorem decoded_policy_regret (sys : System S A O) (choose : ℕ → S → A)
    (estimate : ℕ → S → A → ℚ) (eta tau : ℕ → ℚ)
    (ha : ∀ n s, choose n s ∈ sys.actions s)
    (he : ∀ n s a, a ∈ sys.actions s → |qValue sys n s a - estimate n s a| ≤ eta n)
    (hm : ∀ n s a, a ∈ sys.actions s →
      estimate n s a ≤ estimate n s (choose n s) + tau n)
    (n : ℕ) (s : S) :
    0 ≤ value sys n s - planValue sys s (policyPlan sys choose n s) ∧
    value sys n s - planValue sys s (policyPlan sys choose n s) ≤
      regretBudget (fun k => 2 * eta k + tau k) n
\end{lstlisting}


The conclusion concerns expected finite-horizon value. It is not a pathwise
upper bound on realized loss. This distinction remains in force when the
policy is chosen from calibrated quote intervals later in the theory.

\section{Primitive model certificates and executable planning}
The \path{GuardedModelError} certificate packages an admissible abstraction,
reward discrepancy, transport discrepancy, and table residuals. Its induction
turns those local obligations into a Q-error budget, then applies the guarded
policy regret theorem. The point is to avoid taking a desired whole-policy
bound as an unexplained input.

\paragraph{Statement excerpt.}\path{Komorebi.AVOC.Guarded.model_decoded_policy_regret}. Source: \path{AVOC/GuardedModelError.lean}, line 149.
\begin{lstlisting}[style=AVOCSource]
theorem model_decoded_policy_regret (raw : System S A O) (model : System Z A O)
    (phi : S → Z) (W : ℕ → Z → ℚ) (eps : ℚ) (transport residual : ℕ → ℚ)
    (hc : ModelCertificate raw model phi W eps transport residual)
    (estimate : ℕ → S → A → ℚ) (choose : ℕ → S → A)
    (codeError tau : ℕ → ℚ)
    (ha : ∀ n s, choose n s ∈ raw.actions s)
    (he : ∀ n s a, a ∈ raw.actions s →
      |tableQ model (W n) (phi s) a - estimate n s a| ≤ codeError n)
    (hm : ∀ n s a, a ∈ raw.actions s → estimate n s a ≤ estimate n s (choose n s) + tau n)
    (n : ℕ) (s : S) :
    0 ≤ value raw n s - planValue raw s (policyPlan raw choose n s) ∧
    value raw n s - planValue raw s (policyPlan raw choose n s) ≤
      regretBudget (fun k => 2 * (modelQBudget eps transport residual k + codeError k) + tau k) n
\end{lstlisting}


The executable planner uses finite action lists, bounded depth, and a fallback
selector. \path{choose_allowed} proves admissibility. Its work analysis and
guarantee modules account for truncation and computation charge separately.
The fallback must satisfy its specified legal-action contract. A computation
budget does not justify returning an arbitrary action.

\paragraph{Statement excerpt.}\path{Komorebi.AVOC.ExecutablePlanner.lookahead_policy_regret}. Source: \path{AVOC/PlannerGuarantee.lean}, line 123.
\begin{lstlisting}[style=AVOCSource]
theorem lookahead_policy_regret (m : Model S A O) (sys : Guarded.System S A O)
    (h : Realizes m sys) (fallback : S → A)
    (hf : ∀ s, fallback s ∈ sys.actions s) (R : ℚ)
    (hr : ∀ s a, a ∈ sys.actions s → |sys.reward s a| ≤ R)
    (depth : Nat → Nat) (hd : ∀ k, depth k ≤ k + 1) (n : Nat) (s : S) :
    0 ≤ Guarded.value sys n s - Guarded.planValue sys s
      (Guarded.policyPlan sys (fun k t => choose m fallback (depth k) t) n s) ∧
    Guarded.value sys n s - Guarded.planValue sys s
      (Guarded.policyPlan sys (fun k t => choose m fallback (depth k) t) n s) ≤
      Guarded.regretBudget (fun k => 2 * ((k + 1 - depth k : Nat) : ℚ) * R) n
\end{lstlisting}

\paragraph{Full proof excerpt.}\path{Komorebi.AVOC.ExecutablePlanner.charged_regret}. Source: \path{AVOC/PlannerGuarantee.lean}, line 147.
\begin{lstlisting}[style=AVOCSource]
theorem charged_regret (optimal gross loss computeCost maxCost : ℚ)
    (hlo : 0 ≤ optimal - gross) (hgap : optimal - gross ≤ loss)
    (hc : 0 ≤ computeCost) (hcap : computeCost ≤ maxCost) :
    0 ≤ optimal - chargedReturn gross computeCost ∧
      optimal - chargedReturn gross computeCost ≤ loss + maxCost := by
  unfold chargedReturn
  constructor <;> linarith
\end{lstlisting}


The last declaration makes a simple but consequential distinction. A gross
policy result can be close to optimal while the computation needed to obtain
it is expensive. The charged result adds the permitted computation expense to
the policy loss bound. It is a stated tariff/model quantity, not an empirical
performance measurement of every host operation.

\section{Verification cannot repair a lost specification}
The \path{Verification} module defines a sound gate by implication from gate
acceptance to truth. Its admission theorem follows from that implication. The
\path{Specification} and \path{VerifierEnsemble} modules separate shared
specification error from verifier failure.

\paragraph{Statement excerpt.}\path{Komorebi.AVOC.sound_gate_acceptance}. Source: \path{AVOC/Verification.lean}, line 13.
\begin{lstlisting}[style=AVOCSource]
theorem sound_gate_acceptance
    {C : Type*} {accepted : Finset C} {gate truth : C → Prop}
    (hsound : SoundGate gate truth)
    (hadmit : ∀ c ∈ accepted, gate c) :
    ∀ c ∈ accepted, truth c
\end{lstlisting}

\paragraph{Statement excerpt.}\path{Komorebi.AVOC.cannot_beat_shared_specification_floor}. Source: \path{AVOC/VerifierEnsemble.lean}, line 208.
\begin{lstlisting}[style=AVOCSource]
theorem cannot_beat_shared_specification_floor
    {Ω : Type u} [Fintype Ω]
    (b : RationalBelief Ω)
    {specificationError wrong : Finset Ω}
    (hfloor : specificationError ⊆ wrong) :
    ¬ eventProb b wrong < eventProb b specificationError
\end{lstlisting}


Multiple verifiers can reduce the component of error captured by their joint
failure event under the stated factorization assumptions. They cannot recover
information that the shared specification has erased. The examples include
perfectly correlated failures and genuinely factorized coordinate failures,
so a product-of-marginals conclusion is not used without a premise that
justifies it.

\section{Proof bytes are bound to artifact identity}
The certificate checker is a specified finite derivation language. Its
soundness theorem relates successful checking to the language's propositional
semantics. It is not a universal parser for arbitrary Lean proof files.

\paragraph{Statement excerpt.}\path{Komorebi.AVOC.Certificate.check_sound}. Source: \path{AVOC/CertificateLogic.lean}, line 158.
\begin{lstlisting}[style=AVOCSource]
theorem check_sound (ctx : List Formula) (target : Formula) (proof : Derivation)
    (h : check ctx target proof = true) (v : Nat → Prop) (hc : Satisfies v ctx) :
    Denote v target
\end{lstlisting}


The binding layer records target, environment, dependencies, proof data,
generation, and pending ticket. A successful receipt cannot simply be pasted
onto a different target or a new generation. The exact checks matter because
a valid statement about one artifact is not evidence for an altered artifact.

\paragraph{Statement excerpt.}\path{Komorebi.AVOC.Binding.validated_sound}. Source: \path{AVOC/ArtifactBinding.lean}, line 72.
\begin{lstlisting}[style=AVOCSource]
theorem validated_sound (i : Identity) (t : Ticket) (r : Receipt)
    (h : validate i t r = true) (v : Nat → Prop) (hc : Satisfies v (context i)) :
    Denote v i.statement
\end{lstlisting}

\paragraph{Statement excerpt.}\path{Komorebi.AVOC.Binding.successful_receipt_not_replayable}. Source: \path{AVOC/ArtifactBinding.lean}, line 161.
\begin{lstlisting}[style=AVOCSource]
theorem successful_receipt_not_replayable (s : Slot) (i : Identity) (r : Receipt) (t : Ticket)
    (h : (consume s i r).2 = some t) : (consume (consume s i r).1 i r).2 = none
\end{lstlisting}


One-shot consumption concerns the defined pending-ticket state. It does not
authenticate a remote endpoint. Protected local storage and stable identity
meaning remain application obligations. The source and negative tests include
target, dependency, environment, proof-byte, generation, and stale-ticket
mismatch cases.

\section{Retry control and nonrefundable attempt accounting}
The retry model compares stopping against another trial with a failure
continuation. \path{NoRevival} states that a nonpositive immediate retry value
stays nonpositive after a fault. It is an explicit condition, rather than a
claim about every changing operational environment.

\paragraph{Statement excerpt.}\path{Komorebi.AVOC.Retry.choose_threshold}. Source: \path{AVOC/RetryControl.lean}, line 209.
\begin{lstlisting}[style=AVOCSource]
theorem choose_threshold (h : NoRevival m) (n : Nat) (s : S) :
    choose m (n + 1) s = (m.enabled s && decide (0 < immediate m s))
\end{lstlisting}


The attempt ledger reserves funds and assigns fresh generations. Renewing an
attempt consumes its declared fee and count. Publishing, settling, abandoning,
and replaying have different contracts. A rejected or abandoned attempt does
not erase already incurred charges.

\paragraph{Statement excerpt.}\path{Komorebi.AVOC.Attempts.charge_not_refunded_on_abandon}. Source: \path{AVOC/AttemptAccounting.lean}, line 206.
\begin{lstlisting}[style=AVOCSource]
theorem charge_not_refunded_on_abandon (l r : Ledger) (t : Ticket)
    (h : abandon l t = .ok r) :
    r.charged = l.charged ∧ r.confirmed = l.confirmed ∧ r.terminal = true
\end{lstlisting}

\paragraph{Statement excerpt.}\path{Komorebi.AVOC.Attempts.complete_trace_bounds}. Source: \path{AVOC/AttemptAccounting.lean}, line 321.
\begin{lstlisting}[style=AVOCSource]
theorem complete_trace_bounds (lim : Limits) (es : List Event) (r : Ledger)
    (h : run lim initial es = .ok r) :
    r.confirmed ≤ r.charged ∧ r.charged = (es.map charge).sum ∧
    r.charged ≤ lim.credits ∧ (es.map attemptCount).sum ≤ lim.attempts
\end{lstlisting}


The ledger also rejects stale publication and retry after publication. Even a
zero-fee attempt can exhaust the attempt-count limit. A pure arithmetic credit
bound is distinct from control of the number of external invocations.

\section{Commit recovery is a transition theorem}
Recovery stores a sequence, pre-state, action, committed result, post-state,
and last acknowledged commit. A fresh acknowledgement must agree with the
defined checked transition. An exact duplicate acknowledgement can leave the
checkpoint unchanged. An unfinished intent supplies no durable outcome from
which the machine may invent a report.

\paragraph{Statement excerpt.}\path{Komorebi.AVOC.CommitRecovery.acknowledge_idempotent}. Source: \path{AVOC/CommitRecovery.lean}, line 126.
\begin{lstlisting}[style=AVOCSource]
theorem acknowledge_idempotent (j : Job) (c d : Checkpoint j)
    (r : Commit j) (v : StoreView) (h : acknowledge j c r v = .ok d) :
    acknowledge j d r v = .ok d
\end{lstlisting}

\paragraph{Statement excerpt.}\path{Komorebi.AVOC.CommitRecovery.recovered_safe}. Source: \path{AVOC/CommitRecovery.lean}, line 170.
\begin{lstlisting}[style=AVOCSource]
theorem recovered_safe (j : Job) (c d : Checkpoint j) (o : Outcome j)
    (v : StoreView) (hs : Safe j c) (h : recover j c o v = .ok d) : Safe j d
\end{lstlisting}


These are universal statements about the checkpoint machine. They do not prove
that a physical operation ran exactly once. An ID may be acknowledged once
after a committed result even if an external system executed work in a manner
the model does not observe. A new ID for the same physical work is not excluded
by global ID uniqueness alone.

\subsection{Robust quote recovery has an additional identity}
The later recovery layer retains the whole frozen quote evidence alongside
the selected action. Changed quote evidence invalidates acknowledgement.
The prepared intent is tied to the original robust recommendation. The
acknowledgement transition by itself checks a supplied action and outcome. To
say that action is the original recommendation, the caller must retain the
prepare provenance as well.

\paragraph{Full proof excerpt.}\path{Komorebi.AVOC.Affine.RobustQuoteRecovery.quote_mutation_rejected}. Source: \path{AVOC/RobustQuoteRecovery.lean}, line 127.
\begin{lstlisting}[style=AVOCSource]
@[simp] theorem quote_mutation_rejected (md : QModel n m) (c : Checkpoint md)
    (r : Commit md) (out : RobustQuote.Outcome md) (actual : StoreView md) (ev : EvidenceView n m)
    (h : r.intent.quoteEvidence ≠ ev) :
    acknowledge md c r out actual ev = .error "quote evidence changed after planning" := by
  simp [acknowledge,h]
\end{lstlisting}

\paragraph{Statement excerpt.}\path{Komorebi.AVOC.Affine.RobustQuoteRecovery.acknowledge_checked_transition}. Source: \path{AVOC/RobustQuoteRecovery.lean}, line 228.
\begin{lstlisting}[style=AVOCSource]
theorem acknowledge_checked_transition (md : QModel n m) (c : Checkpoint md)
    (a : Action) (out : RobustQuote.Outcome md) (t : RobustQuote.RuntimeState md)
    (hc : RobustQuote.confirm md c.state a out = some t)
    (hl : c.last ≠ some (commitFor md c a out t)) :
    acknowledge md c (commitFor md c a out t) out (view t) (evidence md) =
      .ok ⟨c.next+1,t,some (commitFor md c a out t)⟩
\end{lstlisting}


The inherited SQLite adapter uses rollback-journal transactions and
\path{synchronous=FULL}. Its intended atomic publication relies on the
SQLite, OS, and filesystem contract. The bounded process-exit suite is
implementation evidence. It is not a Lean proof of disk writes or a universal
power-loss certification.

\section{Verified upgrades retain search costs}
An upgrade authority fixes the protected contract, objective, and core. A
funded trial is charged even when rejected. Promotion requires checked contract
preservation and a strict improvement under the stated criterion. The active
utility can improve while the entire campaign remains unprofitable because
search and switching were expensive.

\paragraph{Statement excerpt.}\path{Komorebi.AVOC.Upgrade.funded_trial_is_charged}. Source: \path{AVOC/VerifiedUpgrade.lean}, line 136.
\begin{lstlisting}[style=AVOCSource]
theorem funded_trial_is_charged [DecidableEq V] (a : Authority V)
    (s : State V) (p : Proposal V) (hf : Funded a s p) :
    (step a s p).next.evaluationSpent = s.evaluationSpent + a.evaluationCost p.candidate
\end{lstlisting}

\paragraph{Statement excerpt.}\path{Komorebi.AVOC.Upgrade.certified_upgrade_campaign}. Source: \path{AVOC/VerifiedUpgrade.lean}, line 260.
\begin{lstlisting}[style=AVOCSource]
theorem certified_upgrade_campaign [DecidableEq V] (a : Authority V) (v : V)
    (hc : a.Contract v) (ps : List (Proposal V)) :
    let s := run a (initial v) ps
    a.Contract s.active ∧
    s.evaluationSpent + s.switchSpent ≤ a.credits ∧
    (s.guaranteedGain : Int) - (s.evaluationSpent : Int) ≤ campaignNet a s
\end{lstlisting}


The campaign theorem packages preserved contract, credit bounds, accounting,
and a net lower bound. Its result is relative to the model's measurements and
tariffs. It gives no claim that arbitrary LLM-generated improvements meet the
checker or that the quoted cost equals physical wall time.

\section{A checked Boolean semantic cost library}
The Boolean backend has three Boolean inputs, hence eight input assignments and
256 Boolean truth signatures. Its program language uses literals, variables,
negation, conjunction, disjunction, and exclusive-or. The cost model is eager
tree nodes. A checked bank supplies one witness program for each signature and
constructor inequalities sufficient for a global lower bound.

\paragraph{Full proof excerpt.}\path{Komorebi.AVOC.BoolClosure.query_globally_minimum}. Source: \path{AVOC/BooleanCostClosure.lean}, line 234.
\begin{lstlisting}[style=AVOCSource]
theorem query_globally_minimum (b : CheckedBank) (p q : Program) (h : Equivalent q p) :
    nodes (query b p) ≤ nodes q := by
  have parts := check_bank_parts b.rows b.valid
  exact witness_global_minimum algebra (cost b.rows) (witness b.rows) parts.2.2 parts.2.1
    (sem p) q ((sem_eq_iff q p).mpr h)
\end{lstlisting}

\paragraph{Statement excerpt.}\path{Komorebi.AVOC.BoolClosure.query_idempotent_cost}. Source: \path{AVOC/BooleanCostClosure.lean}, line 249.
\begin{lstlisting}[style=AVOCSource]
theorem query_idempotent_cost (b : CheckedBank) (p : Program) :
    nodes (query b (query b p)) = nodes (query b p)
\end{lstlisting}


The global minimum quantifies over programs in this declared language and cost
model. It is not universal program minimization. A bank with correct semantic
witnesses but false lower inequalities must be rejected. A checksum protects
the bytes used for a run, but changing a checksum cannot create a valid lower
bound certificate.

The retained library format identifies the input order, operator language,
cost model, and 256-row correspondence. Reopening performs validation again.
An old manifest's initial validation is historical evidence, not a substitute
for the current checker result.

\section{Typed affine semantics and rank-certified factorization}
Affine programs are indexed by input and output dimensions. Sequential
composition must match dimensions. Normalization produces a linear matrix and
bias vector, and \path{normalize_correct} identifies its action on every input.
Finite probes are complete for this affine language because zero determines
the bias and basis vectors determine the linear part.

\paragraph{Statement excerpt.}\path{Komorebi.AVOC.Affine.finite_probe_complete}. Source: \path{AVOC/TypedAffineProgram.lean}, line 56.
\begin{lstlisting}[style=AVOCSource]
theorem finite_probe_complete (a : ∀ {n m}, Atom n m → Signature R n m)
    {n m} (p q : Program Atom n m)
    (h0 : run a p 0 = run a q 0)
    (hb : ∀ j : Fin n, run a p (basis j) = run a q (basis j)) : Equivalent a p q
\end{lstlisting}

\paragraph{Statement excerpt.}\path{Komorebi.AVOC.Affine.context_congr}. Source: \path{AVOC/TypedAffineProgram.lean}, line 94.
\begin{lstlisting}[style=AVOCSource]
theorem context_congr (a : ∀ {n m}, Atom n m → Signature R n m)
    {i o n m} (c : Context Atom i o n m) {p q : Program Atom i o}
    (h : Equivalent a p q) : Equivalent a (plug c p) (plug c q)
\end{lstlisting}


The language restriction is essential. Agreement on zero and a basis does not
identify an arbitrary nonlinear function. Context congruence permits a checked
equivalent component to be substituted inside a typed larger program.

\subsection{A rank witness proves the minimum factor width}
A two-stage factorization represents the linear matrix $A$ as $VU$ through an
intermediate space. The complete upper witness checks composition of both
affine signatures, including their offsets. The lower witness consists of
matrices $L,H$ satisfying $LAH=I_k$, where $k$ is the proposed width. Rank
inequalities then give $k\le\operatorname{rank}(A)$, while the exact width-$k$
factorization gives $\operatorname{rank}(A)\le k$. Any other exact two-stage
factorization must have width at least $k$. No runtime rank calculation is
trusted by this gate.

\paragraph{Statement excerpt.}\path{Komorebi.AVOC.Affine.Factorization.width_minimal}. Source: \path{AVOC/AffineFactorization.lean}, line 66.
\begin{lstlisting}[style=AVOCSource]
theorem width_minimal {f : Signature R n m} {c : Candidate R n m}
    (h : Valid f c) {k : Nat} (a : Signature R n k) (b : Signature R k m)
    (he : ∀ x, eval b (eval a x) = eval f x) : c.width ≤ k
\end{lstlisting}

\paragraph{Statement excerpt.}\path{Komorebi.AVOC.Affine.Factorization.checked_optimal}. Source: \path{AVOC/AffineFactorization.lean}, line 96.
\begin{lstlisting}[style=AVOCSource]
theorem checked_optimal [DecidableEq R] {f : Signature R n m} {c : Candidate R n m}
    (h : check f c = true) :
    (∀ x, execute (plan c) x = eval f x) ∧
    f.linear.rank = c.width ∧
    ∀ (k : Nat) (a : Signature R n k) (b : Signature R k m),
      (∀ x, eval b (eval a x) = eval f x) →
      c.width ≤ k ∧ factorCost n m c.width ≤ factorCost n m k
\end{lstlisting}


The optimizer compares its declared original, dense, and factor menus. A
minimum factor width does not automatically make setup economical at every
call count. The choice and reuse modules retain setup costs and amortization.
The paid and adaptive rank-probe modules also charge information acquisition.
Finding evidence for a cheap factorization is itself an action.

\section{Protected uncertain quotes become robust plans}
\path{QuoteAuthority} binds quotes to exact protected identities. Natural
setup intervals, rational gross intervals, and probability intervals carry
their ordering and probability bounds as proof fields. A compatible
\path{Realization} supplies one true setup, gross value, and probability within
each retained interval.

The conservative nominal model reserves upper setup, credits lower gross,
and uses lower success probability. The true-system comparator keeps the same
conservative upper-setup feasibility guard. It does not silently give an oracle
more spending capacity by changing that reservation rule.

\paragraph{Statement excerpt.}\path{Komorebi.AVOC.Affine.QuoteAuthority.identity_mismatch_rejected}. Source: \path{AVOC/QuoteAuthority.lean}, line 171.
\begin{lstlisting}[style=AVOCSource]
@[simp] theorem identity_mismatch_rejected {n m : Nat} {k : Kind}
    (a : Authority n m k) (q : RawQuote n m)
    (h : q.identity ≠ a.approved.identity) : validate a q = none
\end{lstlisting}

\paragraph{Statement excerpt.}\path{Komorebi.AVOC.Affine.RobustQuote.certified_dominance}. Source: \path{AVOC/RobustQuotePlanning.lean}, line 506.
\begin{lstlisting}[style=AVOCSource]
theorem certified_dominance (md : QModel n m) (h : Nat) (s : State md)
    (a : Action) (_ha : a ∈ LazyStaged.actions (nominal md) s)
    (hd : ∀ b ∈ LazyStaged.actions (nominal md) s, upperQ md h s b ≤ lowerQ md h s a)
    (r : Realization md) :
    ∀ b ∈ (trueSystem md r).actions s,
      Guarded.qValue (trueSystem md r) h s b ≤ Guarded.qValue (trueSystem md r) h s a
\end{lstlisting}

\paragraph{Statement excerpt.}\path{Komorebi.AVOC.Affine.RobustQuote.robust_policy_regret}. Source: \path{AVOC/RobustQuotePlanning.lean}, line 555.
\begin{lstlisting}[style=AVOCSource]
theorem robust_policy_regret (md : QModel n m) (r : Realization md)
    (h : Nat) (s : State md) :
    0 ≤ Guarded.value (trueSystem md r) h s -
      Guarded.planValue (trueSystem md r) s
        (Guarded.policyPlan (trueSystem md r) (robustPolicy md) h s) ∧
    Guarded.value (trueSystem md r) h s -
      Guarded.planValue (trueSystem md r) s
        (Guarded.policyPlan (trueSystem md r) (robustPolicy md) h s) ≤
      Guarded.regretBudget (fun k => 2 * globalQRadius md k) h
\end{lstlisting}


The preceding \path{certified_dominance} is sound but compares the selected
action with every feasible action, including itself. That condition forces
its local Q radius to vanish, as
\path{CalibrationBridge.self_comparison_forces_zero_radius} proves. It is
therefore too restrictive for most uncertain selected actions.

The useful corrected certificate in \path{CalibrationBridge} checks only
distinct feasible competitors. It treats self-comparison by reflexivity and
permits a positive selected radius. Its executable Boolean checker uses the
same competitor-only condition.

\paragraph{Full proof excerpt.}\path{Komorebi.AVOC.Affine.CalibrationBridge.competitor_dominance_sound}. Source: \path{AVOC/CalibrationBridge.lean}, line 21.
\begin{lstlisting}[style=AVOCSource]
theorem competitor_dominance_sound (md : QModel n m) (h : Nat)
    (s : State md) (a : Action) (hd : CompetitorDominance md h s a)
    (r : Realization md) :
    ∀ b ∈ (trueSystem md r).actions s,
      Guarded.qValue (trueSystem md r) h s b ≤
        Guarded.qValue (trueSystem md r) h s a := by
  intro b hb
  by_cases hba : b = a
  · subst b
    exact le_rfl
  · exact (q_between md r h s b).2.trans
      ((hd.2 b hb hba).trans (q_between md r h s a).1)
\end{lstlisting}

\paragraph{Statement excerpt.}\path{Komorebi.AVOC.Affine.CalibrationBridge.checked_competitor_dominance_optimal}. Source: \path{AVOC/CalibrationBridge.lean}, line 57.
\begin{lstlisting}[style=AVOCSource]
theorem checked_competitor_dominance_optimal (md : QModel n m) (h : Nat)
    (s : State md) (a : Action) (hc : competitorDominanceCheck md h s a = true)
    (r : Realization md) :
    ∀ b ∈ (trueSystem md r).actions s,
      Guarded.qValue (trueSystem md r) h s b ≤
        Guarded.qValue (trueSystem md r) h s a
\end{lstlisting}


Action-specific Q intervals can thus certify dominance over distinct
competitors. Otherwise the robust choice maximizes its lower score and the policy obtains the
accumulated Q-radius regret bound. Quotes remain frozen throughout the policy
episode. Spending records conservative reservations. A refund or an equality
between reserved spending and physical expenditure is not part of this
theorem.

\section{Finite calibration is an exact rational construction}
\label{sec:calibration}
For $n$ Bernoulli reports and $k$ successes, the mass polynomial is
\[
 M_{n,k}(p)=\binom nk p^k(1-p)^{n-k}.
\]
On the cell $[i/m,(i+1)/m]$, the checker uses the upper envelope
\[
 U_{n,k,m,i}=\binom nk\left(\frac{i+1}{m}\right)^k
              \left(1-\frac{i}{m}\right)^{n-k}.
\]
It retains a cell when that envelope is strictly larger than
$\alpha/(n+1)$. The returned band is the hull of retained cells. If none is
retained, it returns $[0,1]$. No floating-point likelihood evaluation or
asymptotic normal approximation is used in this formal construction.

\paragraph{Full proof excerpt.}\path{Komorebi.AVOC.FiniteCalibration.excluded_mass_le}. Source: \path{AVOC/FiniteCalibration.lean}, line 153.
\begin{lstlisting}[style=AVOCSource]
theorem excluded_mass_le (n k m : Nat) (hm : 0 < m) {alpha p : ℚ}
    (hp0 : 0 ≤ p) (hp1 : p ≤ 1) (hout : ¬ Contains n k m alpha p) :
    mass n k p ≤ threshold n alpha := by
  obtain ⟨i,him,hil,hiu⟩ := grid_covers m hm hp0 hp1
  have hnot : i ∉ retained n k m alpha := by
    intro hi
    exact hout (retained_cell_in_hull hi hil hiu)
  have henv : envelope n k m i ≤ threshold n alpha := by
    by_contra h
    apply hnot
    exact mem_filter.mpr ⟨mem_range.mpr him, lt_of_not_ge h⟩
  exact (mass_le_envelope n k m i hp0 hp1 hil hiu).trans henv
\end{lstlisting}


Every rational $p\in[0,1]$ lies in a grid cell. If the returned hull excludes
$p$, that covering cell was not retained. Its envelope is therefore at most
the threshold, and the true count mass is at most the envelope. There are
$n+1$ possible counts, giving the total error allowance.

\paragraph{Full proof excerpt.}\path{Komorebi.AVOC.FiniteCalibration.missMass_le}. Source: \path{AVOC/FiniteCalibration.lean}, line 172.
\begin{lstlisting}[style=AVOCSource]
theorem missMass_le (n m : Nat) (hm : 0 < m) {alpha p : ℚ}
    (ha : 0 ≤ alpha) (hp0 : 0 ≤ p) (hp1 : p ≤ 1) :
    missMass n m alpha p ≤ alpha := by
  unfold missMass
  calc
    _ ≤ ∑ _k ∈ range (n+1), threshold n alpha := by
      apply sum_le_sum
      intro k hk
      split
      · unfold threshold; positivity
      · exact excluded_mass_le n k m hm hp0 hp1 ‹_›
    _ = alpha := by
      simp [threshold]
      field_simp
\end{lstlisting}


The grid restricts the reported endpoints, not the true rational probability.
The source theorem handles endpoint probabilities and sample size zero. Its
coverage premise is a binomial count law, whose connection to observations is
proved separately under a source contract.

\subsection{A concrete formal checkpoint}
For ten successes, ten cells, and allowance $1/10$, the retained cells are
$\{6,7,8,9\}$ and the band is $[3/5,1]$. Ten failures return $[0,2/5]$.
Neither sample claims empirical certainty.

\paragraph{Full proof excerpt.}\path{Komorebi.AVOC.FiniteCalibration.ten_successes_retained}. Source: \path{AVOC/FiniteCalibration.lean}, line 188.
\begin{lstlisting}[style=AVOCSource]
theorem ten_successes_retained : retained 10 10 10 (1/10) = {6,7,8,9} := by
  ext i
  by_cases hi : i < 10
  · interval_cases i <;> norm_num [retained,threshold,envelope,cellLower,cellUpper]
  · have h6 : i ≠ 6 := by omega
    have h7 : i ≠ 7 := by omega
    have h8 : i ≠ 8 := by omega
    have h9 : i ≠ 9 := by omega
    simp [retained,hi,h6,h7,h8,h9]
\end{lstlisting}

\paragraph{Full proof excerpt.}\path{Komorebi.AVOC.FiniteCalibration.ten_successes_lower}. Source: \path{AVOC/FiniteCalibration.lean}, line 198.
\begin{lstlisting}[style=AVOCSource]
theorem ten_successes_lower : lower 10 10 10 (1/10) = 3/5 := by
  norm_num [lower,ten_successes_retained,cellLower,min']
\end{lstlisting}


These are exact Lean declarations. The present publication does not rerun
their kernel checks. It reads their source and retained release evidence.

\section{A prospective schedule permits data-dependent selection}
A finite schedule fixes checkpoint identities, sample sizes, grids, and
nonnegative allowances before the data used by those checkpoints. A finite
union bound controls the event that any checkpoint fails. Selection may then
depend on the data without requiring independent checkpoints.

\paragraph{Full proof excerpt.}\path{Komorebi.AVOC.CalibrationSchedule.scheduled_failure_le}. Source: \path{AVOC/CalibrationSchedule.lean}, line 74.
\begin{lstlisting}[style=AVOCSource]
theorem scheduled_failure_le [Fintype J] (law : Law Ω)
    (bad : J → Ω → Prop) (alpha : J → ℚ)
    (h : ∀ j, eventMass law (bad j) ≤ alpha j) :
    eventMass law (fun ω => ∃ j, bad j ω) ≤ ∑ j, alpha j := by
  exact (finite_union_le law bad).trans (sum_le_sum (fun j _ => h j))
\end{lstlisting}

\paragraph{Statement excerpt.}\path{Komorebi.AVOC.CalibrationSchedule.selected_failure_le}. Source: \path{AVOC/CalibrationSchedule.lean}, line 83.
\begin{lstlisting}[style=AVOCSource]
theorem selected_failure_le [Fintype J] (law : Law Ω)
    (bad : J → Ω → Prop) (alpha : J → ℚ)
    (h : ∀ j, eventMass law (bad j) ≤ alpha j) (select : Ω → J) :
    eventMass law (fun ω => bad (select ω) ω) ≤ ∑ j, alpha j
\end{lstlisting}


The price of selection is the sum of the allocated errors across the declared
schedule. The selected checkpoint does not inherit just its own nominal error
without further reasoning. The archive's fair-Boolean adaptive-selection
counterexample makes this distinction executable and formal.

The schedule is not unrestricted optional stopping. Arbitrarily many
unallocated future checkpoints cannot be introduced after seeing the data
and covered by the same finite sum. Within this model, checkpoint dependence
is allowed, but the indexed family and its allowances must be in place.

\section{The primitive source law determines binomial prefix counts}
\path{ObservationStream} uses a finite rational probability space. A sample is
a subset of $\operatorname{Fin}(n)$, marking its successful coordinates.
\path{ConditionalLaw} states that, for any previously specified pattern and
any coordinate not yet specified, the joint mass of that pattern and a fresh
success equals $p$ times the pattern's mass. This avoids dividing by a
conditioning event that might have zero mass.

\paragraph{Definition excerpt.}\path{AVOC/ObservationStream.lean}, line 81.
\begin{lstlisting}[style=AVOCSource]
def ConditionalLaw (law : Law Ω) (n : Nat) (p : ℚ) (stream : Ω → Sample n) : Prop :=
  ∀ (seen target : Sample n) (i : Fin n), i ∉ seen →
    eventMass law (fun ω => Matches seen target (stream ω) ∧ i ∈ stream ω) =
      p * eventMass law (fun ω => Matches seen target (stream ω))
\end{lstlisting}

\paragraph{Statement excerpt.}\path{Komorebi.AVOC.ObservationStream.conditional_implies_product}. Source: \path{AVOC/ObservationStream.lean}, line 132.
\begin{lstlisting}[style=AVOCSource]
theorem conditional_implies_product (law : Law Ω) (n : Nat) (p : ℚ)
    (stream : Ω → Sample n) (h : ConditionalLaw law n p stream) :
    ProductLaw law n p stream
\end{lstlisting}

\paragraph{Statement excerpt.}\path{Komorebi.AVOC.ObservationStream.prefix_count_fiber_eq_binomial}. Source: \path{AVOC/ObservationStream.lean}, line 250.
\begin{lstlisting}[style=AVOCSource]
theorem prefix_count_fiber_eq_binomial (law : Law Ω) (N : Nat) (p : ℚ)
    (stream : Ω → Sample N) (h : ConditionalLaw law N p stream)
    (n : Nat) (hn : n ≤ N) (k : Fin (n + 1)) :
    eventMass law (fun ω => streamCount (fun ω => truncateSample n hn (stream ω)) ω = k) =
      FiniteCalibration.mass n k.val p
\end{lstlisting}


Induction on revealed coordinates establishes the product law. Summing over
patterns with $k$ successes gives the binomial count mass. Truncation inherits
the fresh-coordinate law, so fixed first-$n$ prefixes have the required count
law. Merely checking every one-coordinate marginal is insufficient. The
correlated-pair examples have fair marginals and fail the product law.

The finite rational source contract is a premise. It is not inferred from the
syntax of a log, or from the fact that an operation completed successfully.
That obligation is retained through the final recorder-policy composition.

\section{The raw recorder separates pending intents from terminal records}
\label{sec:raw}
The immutable inherited calibration log cannot turn an unfinished record into
a changed terminal record with the same operation ID. The new recorder fixes
the lifecycle representation by keeping pending intents in a separate list.
Its terminal log continues to obey immutable ID reuse.

\paragraph{Definition excerpt.}\path{AVOC/ObservationRecorder.lean}, line 9.
\begin{lstlisting}[style=AVOCSource]
structure State (Key : Type) where
  pending : List (Nat × Key)
  log : List (Event Key)
  deriving DecidableEq
\end{lstlisting}

\paragraph{Definition excerpt.}\path{AVOC/ObservationRecorder.lean}, line 37.
\begin{lstlisting}[style=AVOCSource]
def step (s : State Key) : Command Key → Option (State Key)
  | .begin id key =>
      if terminalUsed s id then
        if ∃ e ∈ s.log, e.operation = id ∧ e.key = key then some s else none
      else if pendingUsed s id then
        if (id, key) ∈ s.pending then some s else none
      else some ⟨s.pending ++ [(id, key)], s.log⟩
  | .finish e =>
      if e.outcome = .unfinished then none
      else if e ∈ s.log then some s
      else if terminalUsed s e.operation then none
      else if (e.operation, e.key) ∈ s.pending then
        some ⟨s.pending.filter (fun p => p.1 != e.operation), s.log ++ [e]⟩
      else none
\end{lstlisting}


A fresh begin appends an intent. An exact pending begin is unchanged. A late
begin for an already completed ID and the same identity is unchanged and
cannot reopen that operation. Fresh finish requires the matching pending
intent, removes its ID, and appends one terminal event. Exact terminal replay
is unchanged. A changed identity, changed outcome, missing intent, or
unfinished terminal status is rejected.

\subsection{The inductive invariant}
\path{WellFormed} requires duplicate-free pending IDs, duplicate-free terminal
IDs, disjoint pending and terminal IDs, and no unfinished terminal entry. The
initial state satisfies it. The begin and finish cases preserve it, and
induction extends preservation to every successful finite command list.

\paragraph{Full proof excerpt.}\path{Komorebi.AVOC.ObservationRecorder.run_wellFormed}. Source: \path{AVOC/ObservationRecorder.lean}, line 238.
\begin{lstlisting}[style=AVOCSource]
theorem run_wellFormed {cs : List (Command Key)} {s : State Key}
    (h : run cs = some s) : WellFormed s :=
  runFrom_wellFormed initial_wellFormed h
\end{lstlisting}

\paragraph{Statement excerpt.}\path{Komorebi.AVOC.ObservationRecorder.accepted_terminal_retained_once}. Source: \path{AVOC/ObservationRecorder.lean}, line 395.
\begin{lstlisting}[style=AVOCSource]
theorem accepted_terminal_retained_once {s t u : State Key} (hs : WellFormed s)
    {e : Event Key} {cs : List (Command Key)}
    (hf : step s (.finish e) = some t) (hr : runFrom t cs = some u) :
    e ∈ u.log ∧ (u.log.map Event.operation).count e.operation = 1
\end{lstlisting}


Terminal history grows by suffix append. An accepted terminal event remains
exactly once under every later successful continuation. Success and rejection
contribute one Boolean report each. Faults are retained for lifecycle
accountability but contribute no Boolean sample.

\paragraph{Statement excerpt.}\path{Komorebi.AVOC.ObservationRecorder.fault_preserves_samples}. Source: \path{AVOC/ObservationRecorder.lean}, line 341.
\begin{lstlisting}[style=AVOCSource]
theorem fault_preserves_samples {s t : State Key} {id : Nat} {key : Key}
    (h : step s (.finish ⟨id, key, .fault⟩) = some t) (target : Key) :
    samples target t.log = samples target s.log
\end{lstlisting}

\paragraph{Full proof excerpt.}\path{Komorebi.AVOC.ObservationRecorder.inherited_unfinished_terminal_conflict}. Source: \path{AVOC/ObservationRecorder.lean}, line 359.
\begin{lstlisting}[style=AVOCSource]
theorem inherited_unfinished_terminal_conflict (id : Nat) (key : Key)
    (outcome : Outcome) (ht : outcome ≠ .unfinished) :
    record [⟨id, key, .unfinished⟩] ⟨id, key, outcome⟩ = none := by
  apply conflicting_replay
  · simpa using ht
  · exact ⟨⟨id, key, .unfinished⟩, by simp, rfl⟩
\end{lstlisting}


Global operation-ID uniqueness is across identities. It is an invariant of
this reducer's state. It does not forbid a physical caller from using a fresh
ID for duplicated external work.

\section{The generative recorder proves actual sample correspondence}
\label{sec:generative}
The raw recorder receives an external terminal label. The source-generative
recorder has a different interface: begin, report, and fault. A fresh report
constructs its label from the next latent source Boolean for that identity.
Replayed success or rejection is checked before the source is inspected, so a
replay consumes no new Boolean even after source exhaustion.

\paragraph{Definition excerpt.}\path{AVOC/RecorderReveal.lean}, line 46.
\begin{lstlisting}[style=AVOCSource]
def step (s : State Key) : Command Key → Option (State Key)
  | .begin id key => commit s (.begin id key) none
  | .fault id key => commit s (.finish ⟨id, key, .fault⟩) none
  | .report id key =>
      if ⟨id, key, CalibrationLog.Outcome.success⟩ ∈ s.recorder.log then
        commit s (.finish ⟨id, key, .success⟩) none
      else if ⟨id, key, CalibrationLog.Outcome.rejection⟩ ∈ s.recorder.log then
        commit s (.finish ⟨id, key, .rejection⟩) none
      else match s.reveal.remaining key with
        | [] => none
        | bit :: _ =>
            commit s (.finish ⟨id, key, ObservationReveal.checkedOutcome bit⟩) (some key)
\end{lstlisting}


Its state retains the raw recorder, reveal state, emitted ordinary commands,
and effective reveal trace. The \path{Tracks} invariant connects all four.
This is stronger than separately supplying a next-bit equality for each
accepted external event. It proves that the whole generated command list
projects to the raw recorder and the inherited reveal process.

\paragraph{Full proof excerpt.}\path{Komorebi.AVOC.RecorderReveal.run_refines}. Source: \path{AVOC/RecorderReveal.lean}, line 224.
\begin{lstlisting}[style=AVOCSource]
theorem run_refines (source : Key → List Bool) (commands : List (Command Key))
    {s : State Key} (h : run source commands = some s) :
    ObservationRecorder.run s.emitted = some s.recorder ∧
      ∀ key, CalibrationLog.samples key s.recorder.log =
        (ObservationReveal.run source s.effectiveTrace).reports key := by
  have ht := run_tracks source commands h
  refine ⟨ht.recorder_projection, ?_⟩
  intro key
  rw [ht.samples, ht.reveal_projection]
\end{lstlisting}

\paragraph{Full proof excerpt.}\path{Komorebi.AVOC.RecorderReveal.run_samples_eq_take}. Source: \path{AVOC/RecorderReveal.lean}, line 234.
\begin{lstlisting}[style=AVOCSource]
theorem run_samples_eq_take (source : Key → List Bool) (commands : List (Command Key))
    {s : State Key} (h : run source commands = some s) (key : Key) :
    CalibrationLog.samples key s.recorder.log =
      (source key).take (CalibrationLog.samples key s.recorder.log).length := by
  rw [(run_refines source commands h).2 key]
  exact ObservationReveal.run_reports_eq_take source s.effectiveTrace key
\end{lstlisting}


Writing $L_k$ for the actual terminal sample list at identity $k$ and $S_k$ for
the latent source, the conclusion is
\[
 L_k=\operatorname{take}_{|L_k|}(S_k).
\]
Faults, begins, and exact replays do not consume a source position. A fresh
report at an exhausted source is rejected. The finite proof applies to all
successful command lists, including replay-containing lists.

\subsection{The source is constructed, not authenticated}
The theorem's transition definition constructs the next label. It cannot be
used as a test that an arbitrary physical label happens to equal the next
latent bit. A raw log may be well formed and fully resolved without belonging
to the statistical generative model. The physical correspondence is the main
remaining sampling obligation.

\section{Global gates and the final issued-policy theorem}
\label{sec:final}
The final block runs the complete generative command list. Rejection yields no
counts. Any pending operation anywhere yields no counts, even if an unrelated
selected identity has enough reports. A fully resolved run reads each fixed
first-$n$ checkpoint from its actual terminal log. Both false and true reports
count toward readiness. Only true reports count as successes.

\paragraph{Definition excerpt.}\path{AVOC/RecorderPolicy.lean}, line 12.
\begin{lstlisting}[style=AVOCSource]
def gatedCounts (n : J → Nat) (key : J → K)
    (source : Ω → K → List Bool) (commands : Ω → List (RecorderReveal.Command K))
    (j : J) (ω : Ω) : Option (Fin (n j + 1)) :=
  match RecorderReveal.run (source ω) (commands ω) with
  | none => none
  | some s => if s.recorder.pending = [] then
      ObservationPolicy.availableCount (n j) (CalibrationLog.samples (key j) s.recorder.log)
    else none
\end{lstlisting}

\paragraph{Full proof excerpt.}\path{Komorebi.AVOC.RecorderPolicy.pending_run_no_plan}. Source: \path{AVOC/RecorderPolicy.lean}, line 79.
\begin{lstlisting}[style=AVOCSource]
theorem pending_run_no_plan {r c : Nat} (md : QModel r c)
    (n m : J → Nat) (alpha : J → ℚ) (key : J → K)
    (source : Ω → K → List Bool) (commands : Ω → List (RecorderReveal.Command K))
    (selectWarm selectCold : Ω → J) (ω : Ω) (s : RecorderReveal.State K)
    (h : RecorderReveal.run (source ω) (commands ω) = some s)
    (hp : s.recorder.pending ≠ []) :
    recorderPlan md n m alpha key source commands selectWarm selectCold ω = none := by
  apply Affine.ObservedCalibration.observedPlan_none_warm
  exact pending_run_no_count n key source commands ω s h hp (selectWarm ω)
\end{lstlisting}

\paragraph{Full proof excerpt.}\path{Komorebi.AVOC.RecorderPolicy.gatedCounts_faithful}. Source: \path{AVOC/RecorderPolicy.lean}, line 37.
\begin{lstlisting}[style=AVOCSource]
theorem gatedCounts_faithful
    (n : J → Nat) (key : J → K)
    (source : Ω → K → List Bool) (commands : Ω → List (RecorderReveal.Command K))
    (stream : (j : J) → Ω → Sample (n j))
    (hprefix : ∀ j ω, sampleList (stream j ω) = (source ω (key j)).take (n j)) :
    ObservedCalibration.Faithful n (fun j => streamCount (stream j))
      (gatedCounts n key source commands) := by
  intro j ω k hk
  unfold gatedCounts at hk
  cases hr : RecorderReveal.run (source ω) (commands ω) with
  | none => simp [hr] at hk
  | some s =>
    simp only [hr] at hk
    split at hk
    · have hav := ObservationPolicy.availableCount_value (n j)
        (CalibrationLog.samples (key j) s.recorder.log) k hk
      apply Fin.ext
      change k.val = (stream j ω).card
      rw [hav.2, RecorderReveal.run_samples_eq_take (source ω) (commands ω) hr,
        List.take_take, Nat.min_eq_left hav.1, ← hprefix, sampleList_count_true]
    · simp at hk
\end{lstlisting}


The faithfulness proof does not assume that the observed run was completed or
that its samples are a source prefix. It derives the relevant correspondence
on every branch where a count is returned, using the whole-run refinement.
No-count branches remain no-count branches.

\subsection{The certificate being issued}
\path{EpisodeCertified} states existence of a quote realization with the
specified warm/cold true probabilities and deterministic setup/gross values,
together with the frozen-episode Q, value, and expected regret guarantee.
Deterministic setup/gross interval membership remains supplied by a separate
\path{DeterministicRealization}. Sampling only calibrates the probability
fields.

\paragraph{Definition excerpt.}\path{AVOC/CalibratedPolicy.lean}, line 21.
\begin{lstlisting}[style=AVOCSource]
def EpisodeCertified {r c : Nat} (md : QModel r c) (wp cp : ℚ)
    (ws cs : Nat) (wg cg : ℚ) : Prop :=
  ∃ realization : Realization md,
    realization.warm.probability = wp ∧ realization.cold.probability = cp ∧
      realization.warm.setup = ws ∧ realization.cold.setup = cs ∧
      realization.warm.gross = wg ∧ realization.cold.gross = cg ∧
      FrozenEpisodeGuarantee md realization
\end{lstlisting}


For per-key latent sources satisfying \path{ConditionalLaw}, fixed scheduled
checkpoints, positive grids, and nonnegative allowances, the final declaration
bounds the unconditional event of issuing an uncertified episode.

\paragraph{Full proof excerpt.}\path{Komorebi.AVOC.RecorderPolicy.recorder_issued_failure_le}. Source: \path{AVOC/RecorderPolicy.lean}, line 129.
\begin{lstlisting}[style=AVOCSource]
theorem recorder_issued_failure_le {r c : Nat}
    (law : Law Ω) (md : QModel r c) (d : DeterministicRealization md)
    (N : K → Nat) (p : K → ℚ) (base : (k : K) → Ω → Sample (N k))
    (hlaw : ∀ k, ConditionalLaw law (N k) (p k) (base k))
    (n m : J → Nat) (alpha : J → ℚ) (key : J → K)
    (hn : ∀ j, n j ≤ N (key j))
    (hm : ∀ j, 0 < m j) (ha : ∀ j, 0 ≤ alpha j)
    (hp0 : ∀ k, 0 ≤ p k) (hp1 : ∀ k, p k ≤ 1)
    (commands : Ω → List (RecorderReveal.Command K))
    (selectWarm selectCold : Ω → J) (wp cp : ℚ)
    (hwparam : ∀ ω, p (key (selectWarm ω)) = wp)
    (hcparam : ∀ ω, p (key (selectCold ω)) = cp) :
    eventMass law (fun ω => ∃ planned,
      recorderPlan md n m alpha key (fun ω k => sampleList (base k ω)) commands
        selectWarm selectCold ω = some planned ∧
      ¬ EpisodeCertified planned wp cp d.warmSetup d.coldSetup d.warmGross d.coldGross) ≤
      ∑ j, alpha j := by
  have hg := recorder_episode_coverage law md d N p base hlaw n m alpha key hn hm ha
    hp0 hp1 commands selectWarm selectCold wp cp hwparam hcparam
  have hc := eventMass_add_compl law (fun ω => ∃ planned,
      recorderPlan md n m alpha key (fun ω k => sampleList (base k ω)) commands
        selectWarm selectCold ω = some planned ∧
      ¬ EpisodeCertified planned wp cp d.warmSetup d.coldSetup d.warmGross d.coldGross)
  simp only [not_exists, not_and, not_not] at hc
  linarith
\end{lstlisting}


In mathematical notation, with $A=\sum_j\alpha_j$,
\[
 \Prob\bigl(\text{a plan is issued and is not EpisodeCertified}\bigr)\le A.
\]
The probability is over the original finite rational space. Rejected,
unfinished, and insufficient-sample trajectories remain in it. They issue no
plan. The theorem does not divide by the chance of issuance and does not
assert a posterior probability for the observed interval.

\subsection{Why the composition is substantial}
The final theorem no longer takes successful execution, absence of pending
intents, completion, a count law, log-prefix equality, or effective-trace
equality as extra premises. The machine and its refinement derive those
properties where they are needed. The still-explicit assumptions are the
primitive source law, prospective schedule, selected parameter strata, and
deterministic field membership. A policy need not be issued, and no positive
productivity, termination, or collection-cost lower bound follows.

\section{The indexed Lean executor and the JSON boundary}
The executable reducer records the index of the first rejected command. Its
successful state is proved equal to the raw core run for every finite list of
already-decoded commands.

\paragraph{Full proof excerpt.}\path{Komorebi.AVOC.RecorderCodec.replayFrom_toOption}. Source: \path{AVOC/RecorderCodec.lean}, line 18.
\begin{lstlisting}[style=AVOCSource]
theorem replayFrom_toOption {Key : Type} [DecidableEq Key] (s : State Key)
    (index : Nat) (commands : List (Command Key)) :
    (replayFrom s index commands).toOption = runFrom s commands := by
  induction commands generalizing s index with
  | nil => rfl
  | cons c cs ih =>
    cases h : step s c with
    | none => simp [replayFrom, runFrom, h, Except.toOption]
    | some t => simpa [replayFrom, runFrom, h] using ih t (index + 1)
\end{lstlisting}


This theorem covers the executor, not universal correctness of JSON decoding,
host-language parsing, input-file checks, or Python's separate reducer. Those
layers are exercised by finite interface tests. A malformed request is an
input error; a semantically rejected command list returns no final snapshot.
All commands are parsed before an output snapshot is exposed.

The recorder endpoint accepts integer IDs and keys through 2,147,483,647.
Single requests have at most 4,096 commands. Batches have at most 256 words
and 16,384 commands total. The wrapper accepts regular nonsymlink files of at
most 4 MiB. Exact arities and terminal tags matter. The optional injective
registry exports calibration events only after all pending operations are
resolved. Its evidence identifiers are deterministic labels, not signatures
or external proof certificates.

\begin{lstlisting}[style=LeanCode]
["observation-recorder-v1", [
  ["begin", 0, 1],
  ["finish", 0, 1, "success"],
  ["finish", 0, 1, "success"],
  ["begin", 0, 1]
]]
\end{lstlisting}
The defined result has no pending intent and one success terminal record. The
two later commands are exact replays. It is ready in the lifecycle sense, but
readiness does not assert enough calibration samples or any source law.

\section{Positive checks and substantive negative controls}
The retained positive oracle enumerates all command words of lengths zero
through four over sixteen commands, two IDs, and two keys. That is 69,905
Python/reference words. The actual Lean endpoint is compared on lengths zero
through three, plus twelve focused longer scenarios, for 4,381 words. This
precise family is adapter evidence, not an exhaustive proof over arbitrary
IDs or arbitrarily long lists.

The inherited regressions cover robust quote execution, selected-only
workflows, recovery, lazy staged execution, and calibration. Historical v0.40
observation-process evidence remains historical. Its unchanged input pins are
checked, but it is not relabeled as a fresh v0.41 run.

\subsection{What a negative test establishes}
A false theorem attempted with a proof tactic is a useful control when its
failure is substantive. A syntax error, missing import, startup failure,
timeout, or crash is not mathematical rejection. The retained runner classifies
these outcomes separately and requires the exact current membership.

\paragraph{Negative control.}\path{tests/negative/FalseRecorderUnbegunFinish.lean}. Archived as a substantive rejection.
\begin{lstlisting}[style=AVOCSource]
import AVOC.ObservationRecorder
open Komorebi.AVOC.ObservationRecorder
/-- Must fail: a terminal outcome requires the matching previously recorded intent. -/
example : step (initial : State String) (.finish ⟨7, "warm", .success⟩) =
    some ⟨[], [⟨7, "warm", .success⟩]⟩ := by decide
\end{lstlisting}

\paragraph{Negative control.}\path{tests/negative/FalseRecorderReplayDuplicatesSample.lean}. Archived as a substantive rejection.
\begin{lstlisting}[style=AVOCSource]
import AVOC.ObservationRecorder
open Komorebi.AVOC.ObservationRecorder
open Komorebi.AVOC.CalibrationLog
/-- Must fail: exact terminal replay cannot duplicate the checked statistical sample. -/
example : (run [.begin 7 "warm", .finish ⟨7, "warm", .success⟩,
    .finish ⟨7, "warm", .success⟩]).map (fun s => samples "warm" s.log) =
    some [true, true] := by decide
\end{lstlisting}

\paragraph{Negative control.}\path{tests/negative/FalseRecorderPendingPlan.lean}. Archived as a substantive rejection.
\begin{lstlisting}[style=AVOCSource]
import AVOC.RecorderReveal
open Komorebi.AVOC.RecorderReveal
/-- Must fail: an unresolved begun operation blocks issuance. -/
example :
    (run (fun _ : Unit => [true]) [.begin 0 ()]).map (fun s => issue s (7 : Nat)) =
      some (some 7) := by decide
\end{lstlisting}

\paragraph{Negative control.}\path{tests/negative/FalseRecorderReplayConsumesNext.lean}. Archived as a substantive rejection.
\begin{lstlisting}[style=AVOCSource]
import AVOC.RecorderReveal
open Komorebi.AVOC.RecorderReveal
/-- Must fail: replay of the same report consumes no second source outcome. -/
example :
    (run (fun _ : Unit => [false, true]) [.begin 0 (), .report 0 (), .report 0 ()]).map
      (fun s => Komorebi.AVOC.CalibrationLog.samples () s.recorder.log) =
        some [false, true] := by decide
\end{lstlisting}


These controls challenge the absence of an intent, replay sample duplication,
pending-state issuance, and consuming the next bit on replay. Their complete
logs remain in the archive. The positive universal declarations establish the
defined properties. The negative attempts help detect a missing hypothesis or
an accidentally weakened checker.

\subsection{Resource failures were retained}
Two early bounded recorder runs produced no scientific result. A 2 GiB virtual
memory cap prevented Lean thread creation, and an 8 GiB cap prevented startup
mappings. A later diagnostic measured about 8.56 GiB peak virtual space and
3.66 GiB resident memory. The final archived run removed the address-space
cap, retained wall/CPU/log/core bounds, and recorded the inherited 20 GiB cgroup
limit separately. These changes concern execution feasibility. They do not
strengthen the mathematical claim. The failed manifests and diagnostic
correction remain part of the evidence.

\section{Reproduce the formal checkpoint}
\label{sec:reproduce}
Use Linux x86-64 or WSL, Python 3.10 or later, and the separately retained
portable Lean kit with the exact package revisions in \path{lake-manifest.json}.
The source archive omits the runtime and compiled project/dependency objects.
The portable toolchain is an input to reproduction, not a newly built component
of this publication.

\subsection{Check payload and rebuild local objects}
Run the following from the formal project root. The prior ZIP is needed only
to independently repeat the inherited-source comparison. Preserve each command
result rather than treating a command printed in a document as execution.

\begin{lstlisting}[style=LeanCode]
export AVOC_LEAN_KIT=/absolute/path/Komorebi_Portable_Lean_4.34.0-rc2
python3 tools/check_manifest.py
python3 tools/check_inherited_v041.py /absolute/path/AVOC_Formal_v0.40.0.zip
python3 tools/inventory_v041.py
python3 tools/verify_v041.py --prepare
python3 tools/build_v041_parallel.py --seconds 45 --workers 3 --per-target-seconds 180
\end{lstlisting}

Repeat the build command after each invocation exits, until it reports current
222 of 222, zero remaining, and no failed target. Do not repeat
\path{--prepare} between build slices. It removes local objects. The build lock
prevents overlapping release builds, and resumption checks source, imports,
outputs, and logs. No test using those objects should run while preparation or
rebuilding is underway.

\subsection{Audit exact names and run negative controls}
\begin{lstlisting}[style=LeanCode]
python3 tools/verify_v041.py --audit --seconds 600
python3 tools/run_negative_v041.py --seconds 45 --workers 2 --per-file-seconds 90
python3 tools/source_scan_v041.py
\end{lstlisting}
Repeat the negative command until 135 of 135 substantive rejections are
current. The audit expects exactly 2,268 named reports, with only the three
allowed axioms. A zero-hit source scan is a heuristic check. It remains
separate from compilation and the actual axiom reports.

\subsection{Replay actual executable comparisons}
\begin{lstlisting}[style=LeanCode]
python3 checks/test_observation_recorder.py --lean --output results/v0.41.0/oracle_replay.json
python3 tools/run_inherited_regression_v041.py --case all
python3 tools/release_check_v041.py
python3 tools/package_v041.py --output /chosen/AVOC_Formal_v0.41.0.zip
python3 tools/check_manifest.py
\end{lstlisting}
The new output filename preserves the original canonical oracle result and
bounded-run manifest. The inherited runner writes isolated v0.41 evidence and
relocates its fixture paths. The release gate checks exact membership, current
hashes, compiler provenance, the runtime-bound oracle, baseline data,
independent source/tool reviews, and regressions. Changing reviewed source
requires renewed review, not just an edited digest.

\subsection{Run a single recorder example}
\begin{lstlisting}[style=LeanCode]
python3 tools/observation_recorder.py examples/observation_recorder/complete.json
python3 tools/run_recorder_lean.py examples/observation_recorder/complete.json
python3 tools/observation_recorder.py examples/observation_recorder/complete.json --registry examples/observation_recorder/registry.json
\end{lstlisting}
Each request starts from empty and returns a pure snapshot. These commands
perform no paid operation, retry, durable journal mutation, or policy issuance.
The inherited calibration CLI is a separate endpoint. Connecting the two by a
registry exports compatible syntax, not physical statistical authentication.

\section{The physical sampling boundary is sharp}
\label{sec:physical}
Complete logging and eventual resolution do not by themselves establish an
IID stream in finalization order. Consider two independent fair attempted
outcomes. If the successful outcome finishes first whenever the two differ,
the first logged bit is true unless both attempts failed. Its success
probability is $3/4$, not $1/2$. Both attempts can be fully and honestly logged.
The bias comes from order, not missing data.

The retained exact twenty-attempt example completes every success before any
rejection. If at least ten of the twenty fair outcomes succeed, the first ten
completed reports all succeed. The formal calibration construction then
returns $[3/5,1]$, excluding $1/2$, at nominal allowance $1/10$. That event has
mass
\[
 \frac{308333}{524288}>\frac1{10}.
\]
This is exact finite review evidence, not a newly added Lean theorem and not a
counterexample to the source-generative formal theorem. It refutes a proposed
physical interpretation that identifies completion-ordered reports with
unbiased initiation outcomes without a coupling argument.

\subsection{Censoring and conditional guarantees}
The inherited archive also retains outcome-dependent censoring and
conditioning-on-completion examples. An unconditional bad-issuance bound may
hold while the same nominal bound conditioned on completion fails. A fault
that consumes an attempted bit can bias retained reports. A nonconsuming fault
in the formal model has a different meaning. Applications must specify which
process they actually implement.

The remaining route is a precise physical-attempt-to-report-order model or an
initiation-indexed protocol with proved treatment of pending holes, failures,
replay, and ordering. Authentication, honest outcome checks, complete
provenance, identity stability, and stationarity are additional obligations.
No parser, ID invariant, or finite comparison silently supplies them.

\section{What is closed and what remains open}
The current source closes the pure finite recorder lifecycle, whole-run
refinement to source prefixes, indexed Lean execution correspondence, and
composition with fixed scheduled calibration and frozen-episode policy
certificates. These are strong, precisely scoped results. They remain useful
even though physical coupling is separate.

Open scope includes universal host/parser/storage refinement, authenticated
providers and inputs, multi-episode error/reward/cost accumulation, permitted
within-episode updates, empirical workload measurements, arbitrary real or
infinite sources, drift, and unrestricted stopping. Earlier semantic backends
retain their stated languages and cost models. The archive does not claim
universal program minimization or measured LLM productivity.

The companion's fresh source inventory verifies the 195 non-audit module
membership, the recorded line for all 2,268 named declarations, and the 188
theorem-bearing module digests supplied by the declaration-source receipt. It
also checks the consistency of retained build, axiom, negative-test, and final
release summaries. The full machine-readable result includes every immediate
import, every source hash, and the declaration statement excerpts. It labels
all current kernel and actual Lean runtime executions as not run here.

\appendix
\section{Source map and declaration inventory}
The tables below are generated from the actual checkpoint. The machine-readable
inventory retains full qualified names, source paths, source lines, imports,
and SHA-256 digests. The source-only refresh command is
\begin{lstlisting}[style=LeanCode]
python3 scripts/lean_inventory.py
\end{lstlisting}
Run it from the publication root. Its result is a source/receipt check, not a
Lean build. The formal project's own inventory and kernel audit remain the
authoritative source for elaborated declarations and dependencies.

\small
\begin{longtable}{p{0.43\linewidth}r r p{0.30\linewidth}}
\toprule Module & Lines & Named & Immediate local imports\\
\midrule\endhead
\path{AdaptiveProbeCodec} & 120 & 0 & \path{AdaptiveRankProbe}, \path{AffineCodec}\\
\path{AdaptiveProbeControl} & 273 & 22 & \path{ExecutablePlanner}\\
\path{AdaptiveProbeExamples} & 99 & 10 & \path{AdaptiveProbeControl}\\
\path{AdaptiveRankProbe} & 69 & 4 & \path{AdaptiveProbeControl}, \path{AffineRankProbe}\\
\path{AffineCodec} & 155 & 2 & \path{AffineFusion}\\
\path{AffineExamples} & 68 & 12 & \path{AffineFusion}\\
\path{AffineFactorChoice} & 144 & 16 & \path{AffineFactorization}\\
\path{AffineFactorCodec} & 97 & 1 & \path{AffineFactorChoice}, \path{AffineCodec}\\
\path{AffineFactorExamples} & 85 & 16 & \path{AffineFactorCodec}\\
\path{AffineFactorization} & 136 & 16 & \path{AffineFusion}\\
\path{AffineFusion} & 105 & 12 & \path{TypedAffineProgram}\\
\path{AffineProbeCodec} & 116 & 2 & \path{PaidRankProbe}, \path{AffineReuseCodec}\\
\path{AffineProbeExamples} & 67 & 13 & \path{AffineProbeCodec}\\
\path{AffineRankProbe} & 81 & 9 & \path{AffineSearchGuard}\\
\path{AffineReuse} & 141 & 18 & \path{AffineFactorChoice}\\
\path{AffineReuseCodec} & 235 & 11 & \path{AffineSearchGuard}, \path{AffineFactorCodec}\\
\path{AffineReuseExamples} & 65 & 14 & \path{AffineReuseCodec}, \path{AffineFactorExamples}\\
\path{AffineSearchGuard} & 176 & 21 & \path{AffineReuse}\\
\path{AffineSignature} & 112 & 16 & mathlib / none\\
\path{AffineTerminal} & 202 & 17 & \path{AffineReuseCodec}, \path{MultiLevelEconomics}\\
\path{AffineTerminalCodec} & 149 & 3 & \path{AffineTerminal}, \path{MultiLevelCodec}\\
\path{ApproximateBayesian} & 266 & 8 & \path{BayesianEvidence}, \path{ApproximateStochastic}\\
\path{ApproximateBayesianCompression} & 67 & 2 & \path{ApproximateBayesian}, \path{BayesianCompression}, \path{ApproximateComposition}\\
\path{ApproximateComposition} & 130 & 5 & \path{ApproximateStochastic}\\
\path{ApproximateExamples} & 55 & 7 & \path{ApproximateComposition}\\
\path{ApproximateStochastic} & 443 & 19 & \path{StochasticQuotient}\\
\path{ArchiveExamples} & 129 & 13 & \path{ArchiveReplay}, \path{Retrospective}\\
\path{ArchiveReplay} & 233 & 12 & \path{HistoryState}, \path{OpenWorld}, \path{BayesianEvidence}\\
\path{ArtifactBinding} & 198 & 25 & \path{CertificateLogic}\\
\path{ArtifactExamples} & 79 & 14 & \path{CertificateResearch}\\
\path{ArtifactWorkflow} & 296 & 16 & \path{CertificateResearch}, \path{FilteredObservation}\\
\path{ArtifactWorkflowControl} & 170 & 13 & \path{ArtifactWorkflow}\\
\path{ArtifactWorkflowExamples} & 300 & 34 & \path{ArtifactWorkflowSource}\\
\path{ArtifactWorkflowGuarantee} & 151 & 8 & \path{ArtifactWorkflowControl}, \path{BayesianGuardedBridge}\\
\path{ArtifactWorkflowSource} & 109 & 8 & \path{ArtifactWorkflowGuarantee}, \path{BindingCodec}\\
\path{AtomicCommitModel} & 139 & 10 & \path{CommitRecovery}\\
\path{AttemptAccounting} & 330 & 29 & mathlib / none\\
\path{AttemptCodec} & 71 & 2 & \path{AttemptAccounting}\\
\path{AttemptExamples} & 32 & 9 & \path{AttemptAccounting}\\
\path{AttemptRecovery} & 143 & 10 & \path{AttemptAccounting}, \path{CommitRecovery}\\
\path{Bayesian} & 233 & 10 & \path{StochasticBlackwell}\\
\path{BayesianCompression} & 146 & 6 & \path{OpenWorld}\\
\path{BayesianControl} & 275 & 13 & \path{Bayesian}\\
\path{BayesianEvidence} & 161 & 6 & \path{BayesianControl}\\
\path{BayesianExamples} & 145 & 15 & \path{BayesianEvidence}\\
\path{BayesianGuardedBridge} & 168 & 5 & \path{BayesianObservationTransport}, \path{GuardedModelError}\\
\path{BayesianInstability} & 170 & 13 & \path{ApproximateBayesian}, \path{BayesianExamples}\\
\path{BayesianObservationTransport} & 202 & 11 & \path{ApproximateBayesian}\\
\path{BayesianTransportExamples} & 56 & 7 & \path{BayesianObservationTransport}, \path{BayesianExamples}\\
\path{BeliefValueTables} & 77 & 6 & \path{BayesianObservationTransport}\\
\path{BindingCodec} & 109 & 3 & \path{ArtifactBinding}, \path{CertificateCodec}\\
\path{Blackwell} & 141 & 5 & \path{Core}\\
\path{BooleanCostClosure} & 261 & 20 & \path{SemanticCostCertificate}\\
\path{BooleanProgram} & 230 & 23 & mathlib / none\\
\path{BoundedRecommendation} & 105 & 8 & \path{PlannerGuarantee}\\
\path{BoundedSynthesis} & 211 & 24 & \path{BooleanProgram}\\
\path{CalibratedPolicy} & 78 & 2 & \path{CalibrationSchedule}, \path{CalibrationBridge}\\
\path{CalibrationBridge} & 209 & 17 & \path{RobustQuoteCodec}\\
\path{CalibrationCodec} & 108 & 0 & \path{FiniteCalibration}, \path{CalibrationBridge}\\
\path{CalibrationLog} & 112 & 14 & mathlib / none\\
\path{CalibrationSchedule} & 290 & 20 & \path{FiniteCalibration}\\
\path{CertificateCodec} & 64 & 1 & \path{CertificateLogic}\\
\path{CertificateLogic} & 209 & 10 & mathlib / none\\
\path{CertificateResearch} & 143 & 9 & \path{ArtifactBinding}, \path{ResearchPolicy}\\
\path{CertifiedOperationalCatalogue} & 181 & 13 & \path{MultiLevelOperational}\\
\path{CommitRecovery} & 219 & 17 & \path{LiveExecution}\\
\path{Composition} & 37 & 3 & \path{Core}\\
\path{Context} & 153 & 7 & \path{Core}\\
\path{Control} & 222 & 9 & \path{Core}\\
\path{ControlSufficiency} & 395 & 20 & \path{Dynamic}\\
\path{ControlSufficiencyExamples} & 55 & 4 & \path{ControlSufficiency}\\
\path{Core} & 281 & 24 & mathlib / none\\
\path{CostCertificateExamples} & 54 & 8 & \path{SemanticLibrary}\\
\path{DurableUpgrade} & 375 & 22 & \path{VerifiedUpgrade}\\
\path{DurableUpgradeExamples} & 61 & 11 & \path{DurableUpgradeRuntime}\\
\path{DurableUpgradeLaws} & 215 & 15 & \path{DurableUpgrade}\\
\path{DurableUpgradeRuntime} & 140 & 4 & \path{DurableUpgradeLaws}, \path{UpgradeRuntime}\\
\path{Dynamic} & 73 & 4 & \path{Control}\\
\path{EpistemicState} & 77 & 5 & \path{State}\\
\path{ExactRateDistortion} & 495 & 30 & \path{RateDistortionLower}\\
\path{ExactRateDistortionExamples} & 90 & 7 & \path{GreedyRateDistortion}, \path{RateDistortionExamples}\\
\path{Examples} & 99 & 0 & mathlib / none\\
\path{ExecutablePlanner} & 179 & 11 & \path{GuardedModelError}\\
\path{FilteredObservation} & 91 & 8 & \path{BayesianObservationTransport}\\
\path{FiniteCalibration} & 231 & 23 & mathlib / none\\
\path{FiniteEvents} & 82 & 8 & \path{Bayesian}\\
\path{GreedyRateDistortion} & 269 & 18 & \path{ExactRateDistortion}\\
\path{GuardedControl} & 204 & 17 & \path{BayesianControl}\\
\path{GuardedModelError} & 168 & 11 & \path{GuardedPolicyRegret}\\
\path{GuardedPolicyRegret} & 129 & 9 & \path{GuardedControl}\\
\path{HistoryArchive} & 192 & 7 & \path{ArchiveReplay}, \path{HistoryState}\\
\path{HistoryExamples} & 66 & 4 & \path{HistoryState}, \path{OpenWorld}\\
\path{HistoryState} & 265 & 8 & \path{BayesianControl}\\
\path{HorizonBehavior} & 240 & 13 & \path{StateDependentControl}\\
\path{HorizonContext} & 94 & 4 & \path{HorizonBehavior}, \path{Context}\\
\path{HorizonDecisionRateDistortion} & 132 & 6 & \path{QRateDistortionFinite}\\
\path{HorizonExamples} & 36 & 3 & \path{HorizonContext}\\
\path{Investment} & 50 & 4 & mathlib / none\\
\path{LazyStagedCodec} & 252 & 0 & \path{LazyStagedRecovery}, \path{StagedAffineCodec}\\
\path{LazyStagedControl} & 463 & 32 & \path{StagedAffineControl}, \path{StagedAffineOperational}\\
\path{LazyStagedExamples} & 158 & 10 & \path{LazyStagedRecovery}, \path{StagedAffineExamples}\\
\path{LazyStagedRecovery} & 272 & 18 & \path{LazyStagedControl}\\
\path{LiveExecution} & 241 & 19 & \path{PlannerRuntime}\\
\path{LiveExecutionCodec} & 163 & 11 & \path{LiveExecution}\\
\path{ModelBridgeExamples} & 171 & 14 & \path{BayesianGuardedBridge}, \path{BeliefValueTables}\\
\path{MultiLevelCodec} & 161 & 0 & \path{MultiLevelRank}, \path{AffineCodec}, \path{PlannerRuntime}\\
\path{MultiLevelEconomics} & 76 & 6 & \path{MultiLevelRank}, \path{AffineSearchGuard}\\
\path{MultiLevelExamples} & 89 & 7 & \path{MultiLevelInformation}\\
\path{MultiLevelInformation} & 341 & 29 & \path{ExecutablePlanner}\\
\path{MultiLevelOperational} & 187 & 13 & \path{AffineTerminal}, \path{MultiLevelRank}\\
\path{MultiLevelRank} & 204 & 18 & \path{MultiLevelInformation}, \path{AffineRankProbe}\\
\path{ObservationCensoring} & 230 & 22 & \path{CalibrationSchedule}\\
\path{ObservationPolicy} & 166 & 7 & \path{ObservationStream}, \path{ObservationReveal}, \path{ObservedCalibration}\\
\path{ObservationRecorder} & 404 & 37 & \path{CalibrationLog}\\
\path{ObservationReveal} & 186 & 18 & \path{CalibrationLog}\\
\path{ObservationStream} & 296 & 24 & \path{CalibrationSchedule}\\
\path{ObservedCalibration} & 143 & 6 & \path{CalibratedPolicy}\\
\path{OpenWorld} & 146 & 10 & \path{BayesianEvidence}, \path{ApproximateComposition}\\
\path{OperationalExamples} & 169 & 14 & \path{CertifiedOperationalCatalogue}\\
\path{PaidRankProbe} & 245 & 20 & \path{AffineRankProbe}\\
\path{ParallelSchedule} & 110 & 8 & \path{ScheduleCompletion}\\
\path{Pareto} & 169 & 8 & \path{Control}\\
\path{PlannerExamples} & 52 & 11 & \path{WorkflowPlanner}, \path{ArtifactWorkflowExamples}\\
\path{PlannerGuarantee} & 155 & 9 & \path{PlannerWork}\\
\path{PlannerRuntime} & 242 & 11 & \path{WorkflowPlanner}, \path{BoundedRecommendation}\\
\path{PlannerWork} & 132 & 9 & \path{ExecutablePlanner}\\
\path{ProjectedStochasticSufficiency} & 274 & 13 & \path{StochasticSufficiency}\\
\path{ProofScheduling} & 217 & 23 & mathlib / none\\
\path{QRateDistortion} & 289 & 12 & \path{ExactRateDistortion}\\
\path{QRateDistortionExamples} & 85 & 8 & \path{QRateDistortionFinite}\\
\path{QRateDistortionFinite} & 283 & 20 & \path{QRateDistortion}\\
\path{QValueBridge} & 85 & 5 & \path{QRateDistortionFinite}\\
\path{QuoteAuthority} & 292 & 30 & \path{LazyStagedControl}\\
\path{RateDistortionExamples} & 30 & 2 & \path{RateDistortionLower}\\
\path{RateDistortionLower} & 165 & 7 & \path{ApproximateStochastic}\\
\path{RecorderCodec} & 81 & 1 & \path{ObservationRecorder}, \path{CalibrationCodec}\\
\path{RecorderPolicy} & 155 & 7 & \path{RecorderReveal}, \path{ObservationPolicy}\\
\path{RecorderReveal} & 320 & 24 & \path{ObservationRecorder}, \path{ObservationReveal}\\
\path{RecoveryCodec} & 84 & 4 & \path{AtomicCommitModel}, \path{LiveExecutionCodec}\\
\path{RecoveryExamples} & 45 & 6 & \path{RecoveryCodec}\\
\path{RecursiveSynthesis} & 161 & 19 & \path{SemanticIndex}\\
\path{Refinement} & 158 & 9 & \path{Core}\\
\path{RepresentationLattice} & 83 & 7 & \path{OpenWorld}\\
\path{ResearchController} & 266 & 18 & \path{GuardedControl}, \path{SchedulingCertificates}\\
\path{ResearchControllerExamples} & 296 & 29 & \path{ResearchPolicy}, \path{BayesianExamples}\\
\path{ResearchModelGuarantee} & 45 & 1 & \path{BayesianGuardedBridge}, \path{ResearchPolicy}\\
\path{ResearchPolicy} & 156 & 13 & \path{ResearchController}, \path{GuardedPolicyRegret}\\
\path{Retrospective} & 170 & 16 & \path{BayesianCompression}, \path{BayesianExamples}\\
\path{RetryBayesian} & 167 & 16 & \path{RetryControl}, \path{BayesianEvidence}\\
\path{RetryControl} & 217 & 15 & \path{GuardedPolicyRegret}\\
\path{RetryDecisionCodec} & 138 & 3 & \path{RetryLedger}, \path{AttemptCodec}\\
\path{RetryEvaluation} & 69 & 5 & \path{RetryBayesian}\\
\path{RetryExamples} & 128 & 18 & \path{RetryLedger}\\
\path{RetryLedger} & 96 & 5 & \path{RetryEvaluation}, \path{AttemptRecovery}\\
\path{RobustQuoteCodec} & 511 & 11 & \path{RobustQuoteRecovery}, \path{LazyStagedCodec}, \path{AffineReuseCodec}, \path{AffineFactorCodec}\\
\path{RobustQuoteExamples} & 240 & 15 & \path{RobustQuoteRecovery}, \path{LazyStagedExamples}\\
\path{RobustQuotePlanning} & 637 & 40 & \path{QuoteAuthority}, \path{GuardedModelError}, \path{ApproximateStochastic}, \path{ExecutablePlanner}\\
\path{RobustQuoteRecovery} & 238 & 10 & \path{RobustQuoteRuntime}\\
\path{RobustQuoteRuntime} & 295 & 20 & \path{RobustQuotePlanning}\\
\path{ScheduleCompletion} & 169 & 12 & \path{ProofScheduling}\\
\path{SchedulingCertificates} & 77 & 7 & \path{ScheduleCompletion}\\
\path{SchedulingExamples} & 74 & 9 & \path{ParallelSchedule}, \path{SchedulingCertificates}\\
\path{SemanticCostCertificate} & 93 & 6 & \path{BooleanProgram}\\
\path{SemanticIndex} & 250 & 27 & \path{BoundedSynthesis}\\
\path{SemanticLibrary} & 191 & 18 & \path{BooleanCostClosure}\\
\path{SemanticLibraryCodec} & 132 & 2 & \path{SemanticLibrary}\\
\path{SemanticPromotion} & 49 & 3 & \path{SemanticSearchBudget}, \path{SynthesisPromotion}\\
\path{SemanticSearchBudget} & 146 & 8 & \path{SemanticSynthesisWork}\\
\path{SemanticSynthesisExamples} & 51 & 9 & \path{SemanticSearchBudget}\\
\path{SemanticSynthesisRuntime} & 125 & 0 & \path{SemanticPromotion}\\
\path{SemanticSynthesisWork} & 177 & 17 & \path{RecursiveSynthesis}\\
\path{Specification} & 132 & 8 & \path{Verification}\\
\path{StagedAffineCodec} & 167 & 0 & \path{StagedAffineOperational}, \path{AffineTerminalCodec}, \path{PlannerRuntime}\\
\path{StagedAffineControl} & 310 & 22 & \path{AffineTerminal}, \path{ExecutablePlanner}\\
\path{StagedAffineExamples} & 181 & 13 & \path{StagedAffineOperational}, \path{AffineFactorExamples}, \path{AffineReuseExamples}\\
\path{StagedAffineOperational} & 129 & 13 & \path{StagedAffineControl}\\
\path{State} & 70 & 3 & \path{Core}\\
\path{StateDependentControl} & 261 & 12 & \path{ControlSufficiency}\\
\path{Stochastic} & 112 & 6 & \path{Blackwell}\\
\path{StochasticBlackwell} & 503 & 23 & \path{Stochastic}\\
\path{StochasticExamples} & 28 & 3 & \path{StochasticQuotient}\\
\path{StochasticQuotient} & 242 & 8 & \path{ProjectedStochasticSufficiency}\\
\path{StochasticSufficiency} & 116 & 4 & \path{StochasticBlackwell}\\
\path{SynthesisExamples} & 62 & 14 & \path{SynthesisPromotion}\\
\path{SynthesisPromotion} & 136 & 7 & \path{BoundedSynthesis}, \path{DurableUpgradeLaws}\\
\path{SynthesisRetry} & 43 & 4 & \path{SynthesisExamples}, \path{RetryEvaluation}\\
\path{SynthesisRuntime} & 195 & 2 & \path{SynthesisPromotion}\\
\path{TypedAffineProgram} & 153 & 10 & \path{AffineSignature}\\
\path{UpgradeRuntime} & 114 & 4 & \path{UpgradeSpecialization}\\
\path{UpgradeSpecialization} & 193 & 21 & \path{VerifiedUpgrade}, \path{RetryEvaluation}\\
\path{Verification} & 45 & 3 & \path{State}\\
\path{VerifiedUpgrade} & 321 & 27 & mathlib / none\\
\path{VerifierEnsemble} & 261 & 14 & \path{FiniteEvents}, \path{Specification}\\
\path{VerifierExamples} & 113 & 11 & \path{VerifierEnsemble}, \path{BayesianExamples}\\
\path{WorkflowPlanner} & 171 & 11 & \path{PlannerGuarantee}, \path{ArtifactWorkflowSource}\\
\bottomrule
\end{longtable}
\normalsize

\section{Selected exact theorem names}
The following map gives the exact qualified names of every source excerpt in
this companion. The corresponding namespace is the qualified prefix before
the final declaration name. Source lines refer to v0.41.0, not a generated
publication file.

\small
\begin{longtable}{p{0.79\linewidth}p{0.12\linewidth}}
\toprule Qualified declaration and module & Line\\
\midrule\endhead
\path{Komorebi.AVOC.refines_trans}\newline\path{AVOC/Core.lean} & 32\\
\path{Komorebi.AVOC.collision_obstructs_sufficiency}\newline\path{AVOC/Core.lean} & 87\\
\path{Komorebi.AVOC.sufficient_iff_fiber_refinement}\newline\path{AVOC/Core.lean} & 124\\
\path{Komorebi.AVOC.weakLinkUtility_positive_iff}\newline\path{AVOC/Control.lean} & 174\\
\path{Komorebi.AVOC.stochastic_finite_horizon_optimal_action_exists}\newline\path{AVOC/Stochastic.lean} & 81\\
\path{Komorebi.AVOC.quotientStochasticControlSystem_horizonValue_eq}\newline\path{AVOC/StochasticQuotient.lean} & 151\\
\path{Komorebi.AVOC.quotient_stochasticBellmanOptimalAt_iff}\newline\path{AVOC/StochasticQuotient.lean} & 219\\
\path{Komorebi.AVOC.approxBudget_closed_form}\newline\path{AVOC/ApproximateStochastic.lean} & 264\\
\path{Komorebi.AVOC.approximate_stochastic_horizon_value_bound_closed}\newline\path{AVOC/ApproximateStochastic.lean} & 386\\
\path{Komorebi.AVOC.bayesUpdate_mul_predictive}\newline\path{AVOC/Bayesian.lean} & 92\\
\path{Komorebi.AVOC.bayesianPolicyValue_eq_stochasticDecisionValue}\newline\path{AVOC/Bayesian.lean} & 215\\
\path{Komorebi.AVOC.conditionByWeight_l1_stability_combined}\newline\path{AVOC/ApproximateBayesian.lean} & 227\\
\path{Komorebi.AVOC.exists_qSignatureApproximation_iff_fiberDiameter}\newline\path{AVOC/QRateDistortion.lean} & 125\\
\path{Komorebi.AVOC.q_margin_preserves_unique_optimum}\newline\path{AVOC/QRateDistortion.lean} & 256\\
\path{Komorebi.AVOC.Guarded.safe_optimal_plan_exists}\newline\path{AVOC/GuardedControl.lean} & 177\\
\path{Komorebi.AVOC.Guarded.plan_regret_le}\newline\path{AVOC/GuardedPolicyRegret.lean} & 25\\
\path{Komorebi.AVOC.Guarded.decoded_action_gap}\newline\path{AVOC/GuardedPolicyRegret.lean} & 74\\
\path{Komorebi.AVOC.Guarded.decoded_policy_regret}\newline\path{AVOC/GuardedPolicyRegret.lean} & 102\\
\path{Komorebi.AVOC.Guarded.model_decoded_policy_regret}\newline\path{AVOC/GuardedModelError.lean} & 149\\
\path{Komorebi.AVOC.ExecutablePlanner.lookahead_policy_regret}\newline\path{AVOC/PlannerGuarantee.lean} & 123\\
\path{Komorebi.AVOC.ExecutablePlanner.charged_regret}\newline\path{AVOC/PlannerGuarantee.lean} & 147\\
\path{Komorebi.AVOC.sound_gate_acceptance}\newline\path{AVOC/Verification.lean} & 13\\
\path{Komorebi.AVOC.cannot_beat_shared_specification_floor}\newline\path{AVOC/VerifierEnsemble.lean} & 208\\
\path{Komorebi.AVOC.Certificate.check_sound}\newline\path{AVOC/CertificateLogic.lean} & 158\\
\path{Komorebi.AVOC.Binding.validated_sound}\newline\path{AVOC/ArtifactBinding.lean} & 72\\
\path{Komorebi.AVOC.Binding.successful_receipt_not_replayable}\newline\path{AVOC/ArtifactBinding.lean} & 161\\
\path{Komorebi.AVOC.Retry.choose_threshold}\newline\path{AVOC/RetryControl.lean} & 209\\
\path{Komorebi.AVOC.Attempts.charge_not_refunded_on_abandon}\newline\path{AVOC/AttemptAccounting.lean} & 206\\
\path{Komorebi.AVOC.Attempts.complete_trace_bounds}\newline\path{AVOC/AttemptAccounting.lean} & 321\\
\path{Komorebi.AVOC.CommitRecovery.acknowledge_idempotent}\newline\path{AVOC/CommitRecovery.lean} & 126\\
\path{Komorebi.AVOC.CommitRecovery.recovered_safe}\newline\path{AVOC/CommitRecovery.lean} & 170\\
\path{Komorebi.AVOC.Affine.RobustQuoteRecovery.quote_mutation_rejected}\newline\path{AVOC/RobustQuoteRecovery.lean} & 127\\
\path{Komorebi.AVOC.Affine.RobustQuoteRecovery.acknowledge_checked_transition}\newline\path{AVOC/RobustQuoteRecovery.lean} & 228\\
\path{Komorebi.AVOC.Upgrade.funded_trial_is_charged}\newline\path{AVOC/VerifiedUpgrade.lean} & 136\\
\path{Komorebi.AVOC.Upgrade.certified_upgrade_campaign}\newline\path{AVOC/VerifiedUpgrade.lean} & 260\\
\path{Komorebi.AVOC.BoolClosure.query_globally_minimum}\newline\path{AVOC/BooleanCostClosure.lean} & 234\\
\path{Komorebi.AVOC.BoolClosure.query_idempotent_cost}\newline\path{AVOC/BooleanCostClosure.lean} & 249\\
\path{Komorebi.AVOC.Affine.finite_probe_complete}\newline\path{AVOC/TypedAffineProgram.lean} & 56\\
\path{Komorebi.AVOC.Affine.context_congr}\newline\path{AVOC/TypedAffineProgram.lean} & 94\\
\path{Komorebi.AVOC.Affine.Factorization.width_minimal}\newline\path{AVOC/AffineFactorization.lean} & 66\\
\path{Komorebi.AVOC.Affine.Factorization.checked_optimal}\newline\path{AVOC/AffineFactorization.lean} & 96\\
\path{Komorebi.AVOC.Affine.QuoteAuthority.identity_mismatch_rejected}\newline\path{AVOC/QuoteAuthority.lean} & 171\\
\path{Komorebi.AVOC.Affine.RobustQuote.certified_dominance}\newline\path{AVOC/RobustQuotePlanning.lean} & 506\\
\path{Komorebi.AVOC.Affine.RobustQuote.robust_policy_regret}\newline\path{AVOC/RobustQuotePlanning.lean} & 555\\
\path{Komorebi.AVOC.FiniteCalibration.excluded_mass_le}\newline\path{AVOC/FiniteCalibration.lean} & 153\\
\path{Komorebi.AVOC.FiniteCalibration.missMass_le}\newline\path{AVOC/FiniteCalibration.lean} & 172\\
\path{Komorebi.AVOC.FiniteCalibration.ten_successes_retained}\newline\path{AVOC/FiniteCalibration.lean} & 188\\
\path{Komorebi.AVOC.FiniteCalibration.ten_successes_lower}\newline\path{AVOC/FiniteCalibration.lean} & 198\\
\path{Komorebi.AVOC.CalibrationSchedule.scheduled_failure_le}\newline\path{AVOC/CalibrationSchedule.lean} & 74\\
\path{Komorebi.AVOC.CalibrationSchedule.selected_failure_le}\newline\path{AVOC/CalibrationSchedule.lean} & 83\\
\path{Komorebi.AVOC.ObservationStream.conditional_implies_product}\newline\path{AVOC/ObservationStream.lean} & 132\\
\path{Komorebi.AVOC.ObservationStream.prefix_count_fiber_eq_binomial}\newline\path{AVOC/ObservationStream.lean} & 250\\
\path{Komorebi.AVOC.ObservationRecorder.run_wellFormed}\newline\path{AVOC/ObservationRecorder.lean} & 238\\
\path{Komorebi.AVOC.ObservationRecorder.accepted_terminal_retained_once}\newline\path{AVOC/ObservationRecorder.lean} & 395\\
\path{Komorebi.AVOC.ObservationRecorder.fault_preserves_samples}\newline\path{AVOC/ObservationRecorder.lean} & 341\\
\path{Komorebi.AVOC.ObservationRecorder.inherited_unfinished_terminal_conflict}\newline\path{AVOC/ObservationRecorder.lean} & 359\\
\path{Komorebi.AVOC.RecorderReveal.run_refines}\newline\path{AVOC/RecorderReveal.lean} & 224\\
\path{Komorebi.AVOC.RecorderReveal.run_samples_eq_take}\newline\path{AVOC/RecorderReveal.lean} & 234\\
\path{Komorebi.AVOC.RecorderPolicy.pending_run_no_plan}\newline\path{AVOC/RecorderPolicy.lean} & 79\\
\path{Komorebi.AVOC.RecorderPolicy.gatedCounts_faithful}\newline\path{AVOC/RecorderPolicy.lean} & 37\\
\path{Komorebi.AVOC.RecorderPolicy.recorder_issued_failure_le}\newline\path{AVOC/RecorderPolicy.lean} & 129\\
\path{Komorebi.AVOC.RecorderCodec.replayFrom_toOption}\newline\path{AVOC/RecorderCodec.lean} & 18\\
\path{Komorebi.AVOC.Affine.CalibrationBridge.competitor_dominance_sound}\newline\path{AVOC/CalibrationBridge.lean} & 21\\
\path{Komorebi.AVOC.Affine.CalibrationBridge.checked_competitor_dominance_optimal}\newline\path{AVOC/CalibrationBridge.lean} & 57\\
\bottomrule
\end{longtable}
\normalsize

\section{Retained formal and executable reference files}
\small
\begin{itemize}
\setlength{\itemsep}{2pt}
\setlength{\parsep}{0pt}
\item \path{START_HERE.md}, \path{FORMAL_THEORY.md}, \path{PROOF_BOUNDARY.md}:
current scope, mathematical composition, and explicit open boundaries.
\item \path{THEOREM_INDEX.md}, \path{results/v0.41.0/declaration_inventory.json}:
exact 2,268-name source inventory.
\item \path{AUDIT.json}, \path{results/v0.41.0/axiom_audit_summary.json},
\path{results/v0.41.0/full_axiom_audit.log}: retained axiom coverage.
\item \path{results/v0.41.0/build_progress.json},
\path{results/v0.41.0/release_verification.json}: clean target and release receipts.
\item \path{RECORDER_EXECUTABLE.md}, \path{OBSERVATION_EXECUTABLE.md}:
precise endpoint and bounded observation experiment contracts.
\item \path{results/v0.41.0/negative_all/summary.json} and its per-control logs:
exact substantive-rejection evidence.
\item \path{results/v0.41.0/oracle_run/manifest.json},
\path{results/v0.41.0/oracle_results.json}: runtime-bound new endpoint evidence.
\item \path{results/v0.41.0/review_completion_order.py} and its JSON result:
exact ordering obstruction outside the source-generative premise.
\item \path{EXTERNAL_TRUST.md}, \path{RECOVERY_FORMAT.md}, \path{LIBRARY_FORMAT.md}:
physical storage, journal, and certificate format boundaries.
\item \path{REPRODUCE.md}, \path{lean-toolchain}, \path{lake-manifest.json}:
pinned reproduction inputs and workflow.
\end{itemize}
\normalsize

\end{document}
